\section{Introduction}\label{sec:intro}
Milk is perishable in a precise way. Bacteria that grow slowly at 4\degC{} grow quickly at 15\degC, and every hand-off between dairy, truck, depot, shop and home is a chance for the temperature to drift. Predictive microbiology describes this with equations: a primary model gives the bacterial count against time, a secondary model gives the growth rate against temperature, and a tertiary model applies both to a temperature history \cite{perez2013,tarlak2023,stavropoulou2019}. Mathematically the problem is one of non-isothermal kinetics, familiar from thermal analysis and calorimetry: a rate constant that depends on temperature is integrated along a thermal history that is measured, imperfectly, by an instrument \cite{vyazovkin2011,vyazovkin2015,koga2023}. Ozawa's generalized time, which condenses any thermal history into an equivalent time at a reference condition, is the central tool of that field \cite{ozawa1965,ozawa1986,koga2013,malek2013}, and calorimetry has long been used to follow microbial growth and food spoilage directly as a heat flow \cite{schiraldi2019,fessas2017,wadso2011,haman2019,fricke2020}. This paper brings the two traditions together for one practical device, the cold-chain temperature monitor, and asks a question that neither has answered: \emph{given only a temperature record, how well can an inexpensive calculation tell whether milk has passed a limit, and what stops it doing better?}

\subsection{Motivation}\label{sec:why}
The stakes are partly economic and partly about health. Post-pasteurisation contamination with psychrotrophic bacteria causes roughly half to three-fifths of spoilage of fluid milk in the United States \cite{martin2021}, and Gram-negative psychrotrophs that enter after pasteurisation are the main reason high-temperature short-time milk spoils early \cite{ranieri2009}; warm handling therefore wastes food. Foodborne disease more broadly is estimated to cause about 600 million illnesses and 420\,000 deaths a year, with 175\,000 deaths a year in the World Health Organization (WHO) South-East Asia region \cite{who2015}. Those deaths are dominated by other agents and are not specific to milk; we use the figures only as context and model growth, not disease. The microbiological limits a monitor might report against are process-hygiene criteria rather than safety thresholds: the Food Safety and Standards Authority of India (FSSAI) sets an aerobic plate count of $3\times10^4$ colony-forming units (CFU)\,mL$^{-1}$ for pasteurised milk \cite{fssai2024}, which we adopt as our first limit.

Three developments make the question urgent now. First, temperature loggers have become cheap enough to travel with every crate, so the cold chain is being instrumented faster than its data can be interpreted \cite{badiamelis2018,mercier2017,jedermann2014}. Second, dynamic shelf-life labels, which replace a printed date by a prediction from the logged temperature, are being proposed for dairy and meat \cite{koutsoumanis2005,goransson2018,kim2020}. Third, a recent review of fresh dairy concludes that few studies predict the end of shelf life, that endpoint definitions are inconsistent, that few models are validated beyond their development data, and that probabilities are preferable to single estimates \cite{lima2026}. A monitor that is trusted to move or discard milk must therefore come with an honest statement of how often it is wrong and why. Without one, two opposite mistakes are made: money is spent on precise thermometers that cannot improve the prediction, and alarms are believed that are wrong about a third of the time.

\subsection{State of the art}\label{sec:prior}
\textit{Kinetic models.} Growth models are fitted under constant conditions and then applied to changing ones \cite{perez2013,corradini2005}. The modified Gompertz curve \cite{gompertz1825,zwietering1990} and the Baranyi--Roberts model \cite{baranyi1993,baranyi1994} are the standard primary models; the square-root law \cite{ratkowsky1982}, the cardinal-temperature model \cite{rosso1993} and Arrhenius-type laws are the standard secondary models, and ComBase gathers tens of thousands of growth records to fit them \cite{baranyi2004,combase}. For psychrotrophic \textit{Pseudomonas} in milk the square-root law was fitted to 68 growth curves and tested in industry \cite{neumeyer1995,neumeyer1997a}, and growth under fluctuating temperature has been measured directly \cite{lin2016,zwietering1994}.

\textit{Calorimetric studies of microbial growth.} Isothermal calorimetry measures the heat flow of a growing culture and hence its rate. Riva et al.\ proposed it as a general route to the shelf life of foods \cite{riva2001}, and it has been used to follow microbial spoilage of fresh-cut vegetables during shelf life \cite{haman2019}, to measure the temperature dependence of milk fermentation and carrot-juice spoilage, with an apparent activation energy of 120\,kJ\,mol$^{-1}$ for the latter \cite{wadso2011}, to model the growth of lactic bacteria \cite{fessas2017}, and to shorten the detection time of \textit{Pseudomonas} contamination \cite{fricke2020}; a review summarises its role in food science \cite{schiraldi2019}, and inactivation kinetics of foodborne pathogens has also been treated as a thermal-kinetic problem \cite{juneja2011}. These studies measure growth; they do not ask what a temperature record alone can predict.

\textit{Uncertainty analyses.} The nearest studies are Monte Carlo analyses of shelf-life simulations \cite{nauta2000}. For fluid milk, the dairy group at Cornell University modelled spoilage from Gram-negative post-pasteurisation contamination \cite{lau2022} and from psychrotolerant spore-formers \cite{buehler2018}, and deployed the latter as a supply-chain tool \cite{qian2023}. For frozen food, kinetic uncertainty and temperature variability were propagated through nested Monte Carlo loops \cite{giannakourou2020}. In all of them temperature enters as an uncertain input drawn from a distribution.

\textit{Monitoring.} Time--temperature indicators integrate an Arrhenius rate with a fixed activation energy \cite{taoukis1989}, and Fu et al.\ related them to the growth of \textit{Pseudomonas fragi} in dairy products, noting that prior temperature history lengthened the lag \cite{fu1991}; dynamic shelf-life prediction for fish \cite{koutsoumanis2001} led to the Safety Monitoring and Assurance System \cite{koutsoumanis2005}; field studies compare dynamic and printed shelf life \cite{goransson2018}; mean kinetic temperature estimates remaining shelf life \cite{kim2020}; neural networks estimate product temperature from vehicle data \cite{zou2023}; heat-transfer simulation quantifies the cost of cold-chain failures \cite{muller2026,ndraha2018}; and a recent framework couples a growth model with machine-learning components for milk \cite{sofat2026}. Table~\ref{tab:approaches} sets these approaches side by side.

\begin{table*}[!t]\centering\footnotesize
\caption{Approaches to predicting spoilage or quality from a temperature history, and what each leaves open for a temperature monitor}\label{tab:approaches}
\setlength{\tabcolsep}{3.2pt}
\begin{tabular}{@{}p{2.9cm}p{2.3cm}p{2.4cm}p{2.2cm}p{2.1cm}p{4.2cm}@{}}\toprule\rowcolor{dghead}
Approach & Input & Temperature law & Biological variability & Output & Open question for a monitor\\\midrule
Fixed threshold alarm & logged temperature & none & ignored & alarm & how often the alarm is wrong\\
Degree-hours, time above a threshold & logged temperature & linear or step & ignored & index & how close to a growth model it comes\\
Time--temperature indicator, mean kinetic temperature \cite{taoukis1989,kim2020} & logged or physical integral & Arrhenius, fixed $E_a$ & ignored & equivalent temperature & cost of a wrong $E_a$; biology\\
Dynamic shelf life, SMAS \cite{koutsoumanis2001,koutsoumanis2005} & logged temperature & square-root & nominal & remaining life & loss to biology, sensor, sampling\\
Isothermal calorimetry \cite{wadso2011,haman2019,fricke2020} & heat flow at constant $T$ & measured, Arrhenius & measured per sample & rate, detection time & not available in transit\\
Monte Carlo spoilage models \cite{lau2022,buehler2018,qian2023} & simulated temperature & square-root & sampled & share of units spoiled & what is left once temperature is logged\\
Nested uncertainty \cite{giannakourou2020} & simulated temperature & Arrhenius & sampled & shelf-life distribution & monitor accuracy\\
Product-temperature estimation \cite{zou2023} & vehicle and door data & $Q_{10}$ & ignored & product temperature & effect on ranking of spoiled units\\
Heat-transfer coupling \cite{muller2026} & simulated failures & primary + secondary & nominal & shelf-life loss & monitor accuracy\\
Decision intelligence \cite{sofat2026} & sensor stream & growth model + learning & partly & decision & decomposition of error\\
\textbf{This work} & \textbf{sensed, sampled log} & \textbf{square-root, compared with Arrhenius and $Q_{10}$} & \textbf{sampled, with oracle} & \textbf{AUC loss budget, probability} & --\\
\bottomrule\end{tabular}\end{table*}

\subsection{Closest prior work in detail}\label{sec:closest}
Table~\ref{tab:closest} describes the studies closest to ours in more detail: the data they use, their models, how they treat uncertainty, what they found and what they leave open. Table~\ref{tab:criteria} then scores them on eight criteria that a monitor-centred analysis needs, and Table~\ref{tab:findings} sets their numerical findings beside ours. The distinction is one of viewpoint. The sensitivity analyses ask \emph{which unknown input moves the outcome most}; when temperature is unknown it dominates, as they find. We ask \emph{what is left once temperature is logged}; the answer is mostly biology. The two findings fit together.

\begin{table*}[!t]\centering\footnotesize
\caption{The closest prior studies in detail}\label{tab:closest}
\setlength{\tabcolsep}{3.2pt}
\begin{tabular}{@{}p{2.4cm}p{2.7cm}p{2.9cm}p{2.7cm}p{3.0cm}p{2.4cm}@{}}\toprule\rowcolor{dghead}
Study & Data & Model & Uncertainty treatment & Main result & Left open\\\midrule
Lau et al.\ 2022 \cite{lau2022} & isolates from fluid milk; storage temperatures & Gram-negative growth, Monte Carlo & sampled inputs, sensitivity indices & growth rate and storage temperature most influential & accuracy given a log\\
Buehler et al.\ 2018 \cite{buehler2018} & spore-former growth at 6\degC & primary + secondary per isolate & sampled inputs & maximum growth rate most influential & variable temperature\\
Qian et al.\ 2023 \cite{qian2023} & supply-chain temperatures & spore-former model, digital tool & global sensitivity & home and transport temperature, initial count rank highest & what a measured log leaves\\
Giannakourou \& Taoukis 2020 \cite{giannakourou2020} & frozen food, cold-chain data & Arrhenius quality kinetics & nested Monte Carlo & temperature variability dominates & sensor and sampling\\
G\"oransson et al.\ 2018 \cite{goransson2018} & field logs, several chains & dynamic shelf life & none & large gaps, retail weakest & ground truth\\
Kim et al.\ 2020 \cite{kim2020} & storage trials, chicken and milk & mean kinetic temperature & range from extremes & remaining shelf life estimated & biological spread\\
Zou et al.\ 2023 \cite{zou2023} & delivery vehicle data & neural network, $Q_{10}$ & temperature error & 1\degC{} $\Rightarrow$ about 10\,\% shelf-life error & effect on ranking\\
M\"uller et al.\ 2026 \cite{muller2026} & simulated beef chain & heat transfer + growth & scenarios & failures cost up to 60\,\% & monitor accuracy\\
Haman et al.\ 2019 \cite{haman2019} & fresh-cut vegetables & isothermal calorimetry & replicate samples & heat flow tracks spoilage & prediction from a log\\
\bottomrule\end{tabular}\end{table*}

\begin{table*}[!t]\centering\footnotesize
\caption{Closest prior work against eight criteria for a monitor-centred analysis, based on each paper's published text. \yes{} addressed; \half{} partly addressed; \no{} not addressed. C1 starts from a sensed temperature record; C2 scores predictions against a known ground truth; C3, C4, C5 quantify the loss to unobserved biology, to sensor error and to the sampling interval in common units; C6 compares a kinetic score with simple rules on identical trips with paired statistics; C7 assesses calibration of a probability; C8 releases code and data with automated tests}\label{tab:criteria}
\setlength{\tabcolsep}{5pt}
\begin{tabular}{@{}llcccccccc@{}}\toprule\rowcolor{dghead}
Study & Product, question & C1 & C2 & C3 & C4 & C5 & C6 & C7 & C8\\\midrule
Lau et al.\ 2022 \cite{lau2022} & Fluid milk; share of containers over a limit & \no & \half & \half & \no & \no & \no & \no & \no\\
Buehler et al.\ 2018 \cite{buehler2018} & Fluid milk, spore-formers at 6\degC & \no & \half & \half & \no & \no & \no & \no & \no\\
Qian et al.\ 2023 \cite{qian2023} & Fluid milk, supply-chain digital tool & \no & \half & \half & \no & \no & \no & \no & \half\\
Giannakourou \& Taoukis 2020 \cite{giannakourou2020} & Frozen food, nested uncertainty & \no & \no & \half & \no & \no & \no & \no & \no\\
G\"oransson et al.\ 2018 \cite{goransson2018} & Field chains, dynamic vs printed life & \yes & \no & \no & \no & \no & \no & \no & \no\\
Kim et al.\ 2020 \cite{kim2020} & Chicken, milk; mean kinetic temperature & \yes & \half & \no & \no & \no & \no & \no & \no\\
Zou et al.\ 2023 \cite{zou2023} & Delivery; product-temperature estimate & \yes & \half & \no & \half & \no & \no & \no & \no\\
M\"uller et al.\ 2026 \cite{muller2026} & Beef; cost of cold-chain failures & \no & \no & \no & \no & \no & \no & \no & \no\\
Haman et al.\ 2019 \cite{haman2019} & Vegetables; calorimetric shelf-life monitoring & \no & \half & \no & \no & \no & \no & \no & \no\\
Sofat \& Sodhi 2026 \cite{sofat2026} & Milk; decision-intelligence framework & \yes & \half & \no & \half & \no & \half & \no & \no\\
\textbf{This work} & \textbf{Milk; loss budget of a monitor} & \yes & \yes & \yes & \yes & \yes & \yes & \yes & \yes\\
\bottomrule\end{tabular}\end{table*}

\begin{table*}[!t]\centering\footnotesize
\caption{What earlier studies reported and what this work finds on the same point.}\label{tab:findings}
\begin{tabular}{@{}p{3.2cm}p{6.3cm}p{7.0cm}@{}}\toprule\rowcolor{dghead}
Study & Their finding & Our finding\\\midrule
Lau et al.\ \cite{lau2022}, Buehler et al.\ \cite{buehler2018} & Maximum growth rate and storage temperature most influential, with temperature unknown & With temperature logged, growth-rate spread ($cv$ 0.10) alone contributes sd \AnSdB{} to $\ln\CSL$ and $T_{\min}$ spread \AnSdT; biology costs \NumBaseGapMin--\NumBaseGapMax{} AUC\\
Qian et al.\ \cite{qian2023} & Home and transport temperature and initial spore count rank highest & Temperature enters only through $\int\mu\dd t$, which a log measures; what remains is starting load, lag and strain\\
Zou et al.\ \cite{zou2023} & 1\degC{} error $\Rightarrow$ about 10\,\% shelf-life error ($Q_{10}=3$) & Under the square-root law, 1\degC{} $\Rightarrow$ \AnOneDeg\,\% at 4\degC{} and 11\,\% at 10\degC; yet a DS18B20 costs only about 0.01 AUC, because biology dominates\\
G\"oransson et al.\ \cite{goransson2018} & Dynamic and printed shelf life differ widely & A dynamic score beats a fixed time-above-5\degC{} rule by up to 0.07--0.10 AUC\\
Kim et al.\ \cite{kim2020} & Mean kinetic temperature estimates remaining life & Degree-hours, a first-order version of the same idea, comes within \NumLagGainMax{} AUC of the kinetic score\\
M\"uller et al.\ \cite{muller2026} & Short failures shorten beef shelf life by up to 10\,\% & One hour at 13\degC{} uses 3.8\,\% of milk's shelf life (1.2\,\% at 4\degC); hourly sampling misreads it by at most \AnRiem\\
Lin et al.\ \cite{lin2016} & Two 4-h excursions cut time to $10^6$\,CFU\,mL$^{-1}$ from about 90 to about 55\,h & Our reference model gives \NumLinConst{} and \NumLinExc\,h (error $\NumLinErr$\,\%)\\
Neumeyer \cite{neumeyer1995} & \textit{Pseudomonas} generation time 5.6\,h at 4\degC & Same parameters; natural flora in ComBase grow \NumCbBias{} times slower\\
\bottomrule\end{tabular}\end{table*}

\subsection{An illustrative trip}\label{sec:story}
\begin{figure*}[!t]\centering\includegraphics{story}
\caption{One trip, three fates. (a) A carton travels from dairy to home over three days and meets three warm spells (numbered 1--3). (b) Air temperature (dashed line), milk temperature (solid line) and hourly logger readings (dots); the record is identical for every carton on the truck; the shaded band is 2--8\degC. (c) Bacterial counts in three cartons that differ only in biology the logger cannot see (starting load, lag state, strain), computed with the reference model of \secref{sec:math}: carton A solid, carton B dashed, carton C dash-dotted; the horizontal dashed line is the $3\times10^4$\,CFU\,mL$^{-1}$ limit. (d) Shelf life used, computed from the record with Eq.~\eqref{eq:csl}; the thick segments mark the warm spells. (e) Outcome for each carton at the end of the trip}\label{fig:story}\end{figure*}
Figure~\ref{fig:story} shows the problem as a story, read from top to bottom. Three cartons leave the same dairy on the same morning, ride in the same truck, wait at the same depot, sit on the same shop shelf and go home in the same bag (panel a). On the way the truck's door is left open for 90\,min, the shop shelf is restocked for an hour, and at home the carton is left on the counter for two hours. A logger records all of it, so the three cartons have \emph{exactly the same temperature record} (panel b). Inside the cartons the story differs (panel c): carton~C passes the spoilage limit at 50\,h, carton~B ends close to it, and carton~A ends far below. The monitor converts the record into the cumulative fraction of shelf life used (panel d) and must give all three cartons the same answer, which is late for C and early for A and B (panel e). The difference lies in what the logger cannot see: how many bacteria entered each carton after pasteurisation, how ready they were to grow, and how fast that strain grows in the cold. If the biology varies this much, a more accurate thermometer cannot help much, and if the temperature is already logged, sampling it more often adds little. Both statements are intuitive, but neither has been measured.

\subsection{Significance of the study}\label{sec:important}
A monitor's value is decided by the decisions it changes, and those decisions depend on a ranking: which crates to sell first, which to test and which to discard. The area under the ROC curve measures exactly that ranking ability, so a loss budget in AUC tells a dairy, a regulator or a device designer where an extra rupee of effort buys accuracy. Our results show that the answer is not where it is usually spent. A better sensor or faster sampling buys at most 0.02 AUC, while knowledge of the biology is worth up to 0.40. The analysis also gives the thermal-analysis community a clean case in which the generalized-time construction is exact for a living system, in which the temperature law, square-root or Arrhenius, can be tested by its effect on a downstream decision rather than only by goodness of fit, and in which an isothermal calorimetric measurement acquires a direct, quantified value for a supply chain.

\subsection{Advances over prior work}\label{sec:beyond}
\begin{itemize}
\item \textit{A different question.} Earlier studies rank unknown inputs or estimate the cost of a failure; we measure what is left once temperature is logged, judged against a known ground truth.
\item \textit{One unit for every loss.} Biology, sensor, sampling and rule are compared on the same AUC scale, with an oracle as ceiling and paired bootstrap intervals on identical trips.
\item \textit{An exact kinetic reduction.} We prove that both growth models read the history only through a generalized time, which explains why degree-hours, Arrhenius and $Q_{10}$ integrators come close to the growth model and why a cheap sensor suffices.
\item \textit{A calorimetric bridge.} The same reduction gives a master curve for the heat flow of a growing culture and shows that an isothermal calorimetric detection time measures the biology a log cannot see; we quantify how much of the loss it recovers.
\item \textit{Temperature laws tested on decisions.} Four secondary models are compared on ComBase records, and Arrhenius and $Q_{10}$ integrators are compared with the square-root integrator on 240\,000 trips.
\item \textit{Validation and openness.} Parameters are tied to published data and validated on a fluctuating-temperature experiment and on held-out records; code, tests and the compiled database of \DatNFull{} ComBase milk records are released.
\end{itemize}

\subsection{Objectives}\label{sec:aim}
We ask five research questions (RQ). RQ1: how much accuracy is lost to biology that a temperature record cannot see? RQ2: how much is lost to the sensor (a low-cost DS18B20 digital sensor against a PT1000 platinum resistance thermometer) and to the sampling interval? RQ3: does a growth-model score beat simple rules and fixed-law integrators? RQ4: is a probability better calibrated than a single alarm? RQ5: how much of the loss to biology can a calorimetric measurement of the lot recover? We answer each twice, first with mathematics (\secref{sec:math}) and then with a Monte Carlo benchmark (\secref{sec:method}, \secref{sec:stats} and \secref{sec:results}); the two answers agree. A shorter account of the benchmark, without the kinetic theory, the comparison of temperature laws, the error propagation or the statistical development, was released earlier as version 1 of this preprint \cite{singh2026}.

\section{Theory: non-isothermal growth kinetics}\label{sec:math}
\begin{figure*}[!t]\centering\includegraphics{hierarchy}
\caption{The three model layers of predictive microbiology and the members used or compared in this study. The primary model describes the count at fixed conditions, the secondary model the effect of temperature on the rate, and the tertiary model applies both along a temperature history}\label{fig:hierarchy}\end{figure*}
Predictive microbiology works in three layers (Fig.~\ref{fig:hierarchy}). We first describe the primary and secondary models used and compared, then show that, together, they reduce any thermal history to a single generalized time.

This section is written for two readers at once. Each result is stated as a proposition with its proof. Notation is collected in Table~\ref{tab:notation}.

\begin{table}[!htbp]\centering\footnotesize
\caption{Notation and biological meaning.}\label{tab:notation}
\setlength{\tabcolsep}{3pt}
\begin{tabular}{@{}lp{4.55cm}l@{}}\toprule\rowcolor{dghead}
Symbol & Meaning & Unit\\\midrule
$T(t)$ & product temperature at time $t$ & \textdegree C\\
$N,\ y=\ln N$ & bacterial count and its natural log & CFU\,mL$^{-1}$\\
$\mu(T)$ & maximum specific growth rate & h$^{-1}$\\
$b,\ T_{\min}$ & square-root slope; notional minimum growth temperature & h$^{-1/2}$\,K$^{-1}$, \textdegree C\\
$q,\ h_0$ & physiological state; ``work to be done'' in lag, $h_0=\ln(1+1/q_0)$ & --, --\\
$\lambda$ & lag time at constant $T$, $\lambda=h_0/\mu$ & h\\
$N_0,\ N_{\max},\ N_L$ & starting, maximum and limit counts & CFU\,mL$^{-1}$\\
$A,\ L$ & $\log_{10}(N_{\max}/N_0)$, $\log_{10}(N_L/N_0)$ & log$_{10}$\\
$\SL(T)$ & shelf life at constant $T$ (time to reach $N_L$) & h\\
$\CSL$ & cumulative shelf life used & --\\
$\tau(t)$ & biological time, $\int_0^t\mu(T(s))\dd s$ & --\\
$K$ & shelf-life budget in biological time & --\\
$\Delta t$ & sampling interval & h\\
$\beta,\ \varepsilon,\ \delta_q$ & sensor offset, noise, quantisation step & \textdegree C\\
$H$ & label horizon (printed shelf life) & h\\
$cv$ & coefficient of variation of $b$ between strains & --\\
\bottomrule\end{tabular}\end{table}

\subsection{Primary models}\label{sec:primary}
At constant temperature a bacterial population passes through a lag phase, an exponential phase and a stationary phase. The simplest description is exponential growth, $y(t)=y_0+\mu t$, which ignores both lag and saturation. The logistic model adds saturation, $\dd y/\dd t=\mu(1-\mathrm{e}^{y-y_{\max}})$, and the three-phase linear model approximates the curve by three straight segments \cite{perez2013}. Two sigmoid models are used in this study.

\textit{Modified Gompertz model.} Zwietering et al.\ reparametrised the Gompertz function \cite{gompertz1825} in terms of quantities with a biological meaning \cite{zwietering1990}:
\begin{equation}\label{eq:gompertz}
\log_{10}\frac{N(t)}{N_0}=A\exp\!\left\{-\exp\!\left[\frac{\mu_{10}\,\mathrm{e}}{A}(\lambda-t)+1\right]\right\},
\end{equation}
where $A=\log_{10}(N_{\max}/N_0)$ is the asymptotic increase, $\mu_{10}=\mu/\ln10$ the maximum slope in decimal logarithms, and $\lambda$ the lag. The tangent at the inflection point, $t_i=\lambda+A/(\mu_{10}\mathrm{e})$, has slope $\mu_{10}$ and crosses the starting level at $t=\lambda$, which is what makes the parameters interpretable. Equation~\eqref{eq:gompertz} has a closed-form inverse, which is why it is used for the estimator.

\textit{Baranyi--Roberts model.} The reference model for ``real milk'' is the Baranyi--Roberts model \cite{baranyi1993,baranyi1994}, for $y=\ln N$ and a physiological state $q$ that measures how ready the cells are to grow:
\begin{align}
\frac{\dd y}{\dd t} &= \mu(T(t))\,\frac{q}{1+q}\left(1-\mathrm{e}^{\,y-y_{\max}}\right), \label{eq:bar}\\
\frac{\dd q}{\dd t} &= \mu(T(t))\,q, \qquad q(0)=q_0,\quad h_0=\ln\!\left(1+\tfrac1{q_0}\right).\label{eq:barq}
\end{align}
The factor $q/(1+q)$ is a switch that rises from near 0 (cells adapting) to near 1 (cells growing); the factor $1-\mathrm{e}^{y-y_{\max}}$ slows growth as the population nears its ceiling. At constant temperature the solution is closed-form \cite{baranyi1994}:
\begin{align}
y(t)&=y_0+\mu\mathcal{A}(t)-\ln\!\left(1+\frac{\mathrm{e}^{\mu\mathcal{A}(t)}-1}{\mathrm{e}^{\,y_{\max}-y_0}}\right),\label{eq:barcf}\\
\mathcal{A}(t)&=t+\frac1\mu\ln\!\left(\mathrm{e}^{-\mu t}+\mathrm{e}^{-h_0}-\mathrm{e}^{-\mu t-h_0}\right),\nonumber
\end{align}
and the lag is $\lambda=h_0/\mu$: the product $\lambda\mu=h_0$ is constant for cells with the same history, so $h_0$ is the ``amount of work'' the cells must do before growing, whatever the temperature.


\subsection{Secondary models}\label{sec:secondary}
The square-root law of Ratkowsky et al.\ \cite{ratkowsky1982} states that, below the optimum temperature,
\begin{equation}\label{eq:sqrt}
\sqrt{\mu(T)} = b\,(T - T_{\min}), \qquad T > T_{\min},
\end{equation}
so $\mu(T)=b^2(T-T_{\min})^2$ is a parabola that touches zero at $T_{\min}$. Given $n$ measured rates $\mu_i$ at temperatures $T_i$, we fit Eq.~\eqref{eq:sqrt} by ordinary least squares on $\sqrt{\mu_i}=a+bT_i$:
\begin{equation}\label{eq:ols}
\hat b=\frac{\sum_i (T_i-\bar T)(\sqrt{\mu_i}-\overline{\sqrt{\mu}})}{\sum_i (T_i-\bar T)^2},\qquad \hat T_{\min}=-\frac{\hat a}{\hat b},
\end{equation}
with $\hat a=\overline{\sqrt{\mu}}-\hat b\bar T$, and intervals by resampling records with replacement \cite{efron1993}. The law also fixes the familiar $Q_{10}$, the factor by which growth speeds up for a 10\degC{} rise:
\begin{equation}\label{eq:q10}
Q_{10}(T)=\frac{\mu(T+10)}{\mu(T)}=\left(\frac{T+10-T_{\min}}{T-T_{\min}}\right)^{2}.
\end{equation}
Unlike a constant $Q_{10}$, this ratio is largest in the cold: with $T_{\min}=-7.7$\degC{} \cite{neumeyer1995} it is 3.4 at 4\degC{} and 2.1 at 15\degC. The local sensitivity of the rate to temperature is
\begin{equation}\label{eq:sens}
\frac{\dd\ln\mu}{\dd T}=\frac{2}{T-T_{\min}},
\end{equation}
which is 0.17 per \textdegree C at 4\degC{} and 0.11 at 10\degC.


Five secondary models are common in food microbiology and thermal analysis. The square-root law, Eq.~\eqref{eq:sqrt}, is used throughout; for the whole biokinetic range it is extended to $\sqrt{\mu}=b(T-T_{\min})\{1-\exp[c(T-T_{\max})]\}$ \cite{ratkowsky1982}, which reduces to Eq.~\eqref{eq:sqrt} well below the maximum temperature $T_{\max}$. The Arrhenius law and its three-parameter extension by Davey are
\begin{gather}
\ln\mu=\ln A_0-\frac{E_a}{RT_K},\label{eq:arr}\\
\ln\mu=C_0+\frac{C_1}{T_K}+\frac{C_2}{T_K^2},\label{eq:davey}
\end{gather}
with $T_K$ the absolute temperature. A constant temperature coefficient gives the exponential law
\begin{equation}\label{eq:q10law}
\ln\mu=a+cT,\qquad Q_{10}=\mathrm{e}^{10c},
\end{equation}
and the cardinal-temperature model of Rosso et al.\ describes the rate through $T_{\min}$, the optimum and $T_{\max}$ \cite{rosso1993}. Below the optimum, which is the only range relevant to a cold chain, the cardinal model and the extended square-root law reduce to curves of the same shape as Eq.~\eqref{eq:sqrt}, so we compare the four laws of Eqs.~\eqref{eq:sqrt} and \eqref{eq:arr}--\eqref{eq:q10law}.

Each law was fitted by ordinary least squares, the square-root law on $\sqrt{\mu}$ and the others on $\ln\mu$, to the \NumCbN{} ComBase records of spoilage flora in pasteurised milk and to the 46 \textit{L.~monocytogenes} fitting records (Table~\ref{tab:secondary}, Fig.~\ref{fig:secondary}). The statistics used are defined in \secref{sec:stats}. Within the range of the data all four laws fit about equally well, and the square-root law has the lowest AICc for both organisms. The decisive difference appears in the cold. When each law is fitted to the records at or above 8\degC{} and used to predict the records below 8\degC, the square-root law is almost unbiased (bias factor 1.01 for spoilage flora and 1.15 for \textit{Listeria}), while the Arrhenius and constant-$Q_{10}$ laws overpredict by 16--18\,\% and 46--49\,\%. The curvature of Eq.~\eqref{eq:sqrt} on an Arrhenius plot is exactly what Ratkowsky et al.\ observed \cite{ratkowsky1982}. Fitted over the whole range, Eq.~\eqref{eq:arr} gives $E_a=\SecEaSp$\,kJ\,mol$^{-1}$ for spoilage flora and \SecEaLis\,kJ\,mol$^{-1}$ for \textit{Listeria}, of the same order as the 120\,kJ\,mol$^{-1}$ measured calorimetrically for carrot-juice spoilage \cite{wadso2011}; the corresponding constant $Q_{10}$ values are \SecQtenSp{} and \SecQtenLis. On held-out \textit{Listeria} records all four laws perform alike, so the advantage of the square-root law is specific to extrapolation into the cold, which is where a cold chain operates.

\TabSecondary

\begin{figure*}[!t]\centering\includegraphics{secondary}
\caption{Secondary models fitted to ComBase records. (a) Spoilage flora in pasteurised milk and (b) \textit{L.~monocytogenes}: growth rates (points) with the square-root (solid), Arrhenius (dashed), constant-$Q_{10}$ (dotted) and Arrhenius--Davey (dash-dotted) fits. (c) Accuracy factor for the records below 8\degC{} when each law is fitted at or above 8\degC; plain bars spoilage flora, hatched bars \textit{Listeria}. Data from ComBase \cite{combase}}\label{fig:secondary}\end{figure*}

\subsection{Relation to Arrhenius kinetics and the apparent activation energy}
In thermal analysis the rate constant usually follows Arrhenius, $k(T)=A\exp(-E/RT)$ \cite{arrhenius1889}. Ratkowsky et al.\ found that bacterial growth rates bend away from an Arrhenius line and are described better by Eq.~\eqref{eq:sqrt} \cite{ratkowsky1982}. The two laws can be compared through the apparent activation energy, defined as for any rate,
\begin{equation}\label{eq:ea}
E_{\mathrm{app}}(T)=R\,T_K^{2}\,\frac{\dd\ln\mu}{\dd T}=\frac{2R\,T_K^{2}}{T-T_{\min}},
\end{equation}
with $T_K$ the absolute temperature. For spoilage flora with $T_{\min}=-7.7$\degC{} it falls from 109\,kJ\,mol$^{-1}$ at 4\degC{} to 84 at 8\degC{} and 66 at 13\degC; for \textit{Listeria} ($T_{\min}=-2.3$\degC) from 203 to 89\,kJ\,mol$^{-1}$ (Fig.~\ref{fig:rates}c). An isoconversional analysis of growth data would therefore report an activation energy that varies with temperature, which is the usual signature of a process that is not a single Arrhenius step \cite{vyazovkin2011}. For the monitor this matters because a time--temperature indicator or a mean-kinetic-temperature rule built on one fixed $E$ \cite{taoukis1989,kim2020} treats warm spells differently from the bacteria.

\begin{figure*}[!t]\centering\includegraphics{rates}
\caption{Temperature dependence of growth. (a) Square-root law for spoilage flora (baseline, solid line) and \textit{L.~monocytogenes} (ComBase fit, dashed line). (b) Shelf life at a steady temperature, $\SL=K/\mu$, for the $3\times10^4$\,CFU\,mL$^{-1}$ limit (logarithmic axis). (c) Equivalent Arrhenius activation energy from Eq.~\eqref{eq:ea}}\label{fig:rates}\end{figure*}

\subsection{Reduction to a generalized time}
The central observation of this paper is that the growth equations care about temperature only through one accumulated quantity.

\begin{definition}[Biological time]
For a temperature history $T(\cdot)$, the biological time elapsed by $t$ is $\tau(t)=\int_0^t \mu(T(s))\dd s$.
\end{definition}

\begin{proposition}[Generalized-time reduction]\label{prop:clock}
Let $\mu(T)$ drive both Eqs.~\eqref{eq:bar}--\eqref{eq:barq}. Then $y$ and $q$ depend on the history only through $\tau(t)$: there are functions $Y,Q$ with $y(t)=Y(\tau(t))$, $q(t)=Q(\tau(t))$, where $Y,Q$ solve the same equations with $\mu\equiv1$. In particular $q(t)=q_0\mathrm{e}^{\tau(t)}$, the time at which $N$ reaches $N_L$ is the first $t$ with $\tau(t)\ge\tau^*$, and $\tau^*$ depends only on $(h_0,N_0,N_{\max},N_L)$, not on temperature.
\end{proposition}

The construction is the same as Ozawa's generalized time in thermal analysis, $\theta=\int_0^t\exp(-E_a/RT)\dd s$, which condenses any thermal history into the time an isothermal experiment at infinite temperature would need to reach the same conversion \cite{ozawa1965,ozawa1970,ozawa1986}. In that notation the kinetic equation $\dd\alpha/\dd t=A\exp(-E_a/RT)f(\alpha)$ becomes $\dd\alpha/\dd\theta=Af(\alpha)$, and plots of the rate against conversion in generalized time are master plots that do not depend on the temperature programme \cite{malek2013,koga2013}. Proposition~\ref{prop:clock} is the same statement for bacterial growth: the right-hand side of Eqs.~\eqref{eq:bar}--\eqref{eq:barq} divided by $\mu$ is the kinetic model function, $\tau$ is the generalized time with the square-root law in place of the Arrhenius factor, and the count plotted against $\tau$ is a master curve. Figure~\ref{fig:clock} checks the proposition numerically: two very different histories with the same biological time give different growth paths in elapsed time but exactly the same curve in biological time.

\begin{figure*}[!t]\centering\includegraphics{clock}
\caption{The generalized-time reduction (Proposition~\ref{prop:clock}) shown numerically with the Baranyi--Roberts model. (a) Two thermal histories with the same total biological time: a steady 5.5\degC{} (trip 1, solid line) and 4\degC{} with a 6-h spell at 13\degC{} (trip 2, dashed line). (b) Biological time $\tau$. (c) Bacterial counts against elapsed time. (d) Bacterial counts against biological time; the two curves coincide}\label{fig:clock}\end{figure*}


\subsection{The cumulative shelf-life accumulator}
A logger cannot solve differential equations cheaply, so the edge estimator uses the modified Gompertz curve \cite{zwietering1990}, $\log_{10}(N/N_0)=A\exp\{-\exp[\mu_{10}\mathrm{e}(\lambda-t)/A+1]\}$ with $\mu_{10}=\mu/\ln10$. Solving for the time at which $\log_{10}(N/N_0)=L$ gives
\begin{equation}\label{eq:sl}
\SL(T) = \frac{h_0}{\mu(T)} + \frac{A\ln10}{\mu(T)\,\mathrm{e}}\left[1-\ln\!\left(-\ln\frac{L}{A}\right)\right]=\frac{K}{\mu(T)},
\end{equation}
with the constant
\begin{equation}\label{eq:K}
K = h_0 + \frac{A\ln 10}{\mathrm{e}}\left[1-\ln\!\left(-\ln\frac{L}{A}\right)\right].
\end{equation}
A reading $T_i$ held over $\Delta t_i$ uses a fraction $\Delta t_i/\SL(T_i)$ of the shelf life, and the cumulative shelf life used (CSL) is
\begin{equation}\label{eq:csl}
\CSL = \sum_i \frac{\Delta t_i}{\SL(T_i)}, \qquad \text{limit reached when } \CSL\ge 1.
\end{equation}
This costs one table lookup and one addition per reading. The estimator is deliberately a different primary model from the reference, so the comparison is not circular.

\begin{proposition}[The accumulator is a function of generalized time]\label{prop:csl}
With continuous readings, $\CSL(t)=\tau(t)/K$. Hence the Gompertz accumulator and the Baranyi--Roberts reference are both monotone functions of the same biological time, and they cross their limits at $\tau=K$ and $\tau=\tau^*$ respectively.
\end{proposition}


\begin{worked}
With the baseline parameters ($b=0.030$, $T_{\min}=-7.7$\degC), $\mu(4)=(0.030\times11.7)^2=\AnMuFour$\,h$^{-1}$ and $\mu(13)=\AnMuThirteen$\,h$^{-1}$. For a limit of $3\times10^4$\,CFU\,mL$^{-1}$, $K=\AnKnom$, so the shelf life is $\AnKnom/\AnMuFour\approx83$\,h at a steady 4\degC{} and $\AnKnom/\AnMuThirteen\approx26$\,h at 13\degC. One hour at 4\degC{} consumes 1.2\,\% of the shelf life and one hour at 13\degC{} consumes 3.8\,\%. A 90-minute spell at 13\degC{} therefore costs about 4\,\% of shelf life beyond what the same time at 4\degC{} would have cost; in Fig.~\ref{fig:story}d each warm spell appears as such a step in the gauge.
\end{worked}

\subsection{Degree-hours as a first-order approximation}
\begin{corollary}[Degree-hours as a first-order approximation]\label{cor:dh}
Expand $\mu$ about a reference $T_r$. Under Eq.~\eqref{eq:sqrt},
\begin{multline}\label{eq:dh}
\CSL=\frac{b^2}{K}\Big[(T_r-T_{\min})^2\,t+2(T_r-T_{\min})\!\int_0^t\!(T-T_r)\dd s\\+\int_0^t(T-T_r)^2\dd s\Big].
\end{multline}
\end{corollary}
For trips of equal length the first term is common to all trips, the second is a signed degree-hour count about $T_r$, and the third gives extra emphasis to large excursions. With $T_r=4$\degC, an excursion to 13\degC{} gets $(9)^2=81$ from the quadratic term against $2\times11.7\times9=211$ from the linear term, so degree-hours capture about 72\,\% of its cost; for an excursion to 30.7\degC{} they capture only 47\,\%. A degree-hour count should therefore rank trips almost as well as $\CSL$, losing a little on hot excursions, which is what the simulation finds (\secref{sec:results}).


\subsection{Fixed-law integrators as generalized times}\label{sec:mkt}
Every single-parameter temperature integrator is a generalized time with some rate law. A time--temperature indicator or mean kinetic temperature (MKT) uses the Arrhenius factor with a fixed $E_a$ \cite{taoukis1989,kim2020},
\begin{gather}
\theta_E=\int_0^t\exp\!\left[-\frac{E_a}{R}\left(\frac{1}{T_K}-\frac{1}{T_{\mathrm{ref}}}\right)\right]\dd s,\label{eq:mkt}\\
T_{\mathrm{MKT}}=\frac{E_a/R}{E_a/(RT_{\mathrm{ref}})-\ln(\theta_E/t)},\nonumber
\end{gather}
a $Q_{10}$ rule uses $Q_{10}^{(T-T_{\mathrm{ref}})/10}$, and degree-hours use $(T-T_{\mathrm{ref}})^+$. For trips of equal duration $T_{\mathrm{MKT}}$ is a monotone function of $\theta_E$, so it ranks trips in the same order. Any of these integrators therefore differs from $\CSL$ only in the weight it gives an hour at temperature $T$ relative to an hour at the reference. Table~\ref{tab:weights} lists these weights. Around a warm spell to 13\degC{} the laws differ by less than a factor of two, but for a 30\degC{} excursion the Arrhenius weights are three to nine times the square-root weight. An Arrhenius integrator therefore over-penalises hot, short spells relative to the growth they cause, and the larger its $E_a$ the more it does so; this is the mechanism behind the benchmark result in \secref{sec:results}.

\begin{table}[!htbp]\centering\footnotesize
\caption{Weight of one hour at temperature $T$ relative to one hour at 4\degC{} under each temperature law (spoilage flora; square-root law with $T_{\min}=-7.7$\degC)}\label{tab:weights}
\setlength{\tabcolsep}{4pt}
\begin{tabular}{@{}lcccc@{}}\toprule\rowcolor{dghead}
$T$/\textdegree C & Square-root & \shortstack{Arrhenius\\\SecEaSp\,kJ\,mol$^{-1}$} & \shortstack{Arrhenius\\120\,kJ\,mol$^{-1}$} & $Q_{10}=3$\\\midrule
8 & 1.80 & 1.77 & 2.10 & 1.55\\
13 & 3.13 & 3.52 & 5.14 & 2.69\\
17.5 & 4.64 & 6.43 & 11.2 & 4.41\\
30.7 & 10.8 & 33.8 & 97.2 & 18.8\\
\bottomrule\end{tabular}\end{table}

\subsection{Connection with isothermal calorimetry}\label{sec:calotheory}
Isothermal calorimetry follows microbial growth through the heat flow $P$ of a sample \cite{wadso2011,fessas2017,schiraldi2019}. During growth the heat flow is taken as proportional to the rate at which cells are formed, $P=Y_Q V\,\dd N/\dd t$, where $Y_Q$ is the heat released per cell formed and $V$ the sample volume \cite{fessas2017,fricke2020}. Under the Baranyi--Roberts model, $\dd N/\dd t=N\,\dd y/\dd t=\mu(T)\,N\,\frac{q}{1+q}(1-N/N_{\max})$, and Proposition~\ref{prop:clock} turns this into a statement about calorimetric curves.

\begin{proposition}[Calorimetric master curve]\label{prop:calo}
Let $Y_Q$ be constant and let the growth follow Eqs.~\eqref{eq:bar}--\eqref{eq:barq}. Then, for any temperature programme,
\begin{equation}\label{eq:calo}
\frac{P(t)}{\mu\big(T(t)\big)}=Y_QV\,G\big(\tau(t)\big),\qquad Q(t)=\int_0^tP\dd s=Y_QV\big[N(\tau(t))-N_0\big],
\end{equation}
where $G(\tau)=N\frac{q}{1+q}(1-N/N_{\max})$ is evaluated on the solution with $\mu\equiv1$. The calorimetric conversion $\alpha=Q/Q_\infty=(N-N_0)/(N_{\max}-N_0)$ therefore obeys $\dd\alpha/\dd\tau=f(\alpha)$ with a model function $f$ that does not depend on the temperature programme.
\end{proposition}

Equation~\eqref{eq:calo} is the biological counterpart of the master plots of thermal analysis, $\dd\alpha/\dd\theta=Af(\alpha)$ \cite{malek2013,koga2013}: thermograms recorded under different programmes differ in elapsed time, but divided by $\mu(T)$ and plotted against $\tau$ they coincide (Fig.~\ref{fig:calorimetry}a,b). Two uses follow. First, a single non-isothermal calorimetric run, for example a stepwise programme, can be used to test the square-root generalized time and to estimate $b$ and $T_{\min}$ for a given product, in the same way as Ozawa's method uses runs at several heating rates \cite{ozawa1965,koga2013}; a calorimeter that measures at several temperatures at once \cite{wadso2011} gives $\mu(T)$ and hence $E_{\mathrm{app}}$, Eq.~\eqref{eq:ea}, directly.

Second, calorimetry measures the biology that a temperature log cannot. Consider a sample of the lot incubated isothermally at $T_{\mathrm{ref}}$ until the heat flow exceeds a detection threshold, the detection-time method used to screen for contamination \cite{fricke2020}. If detection corresponds to a count $N_{\mathrm{det}}$, Proposition~\ref{prop:clock} gives
\begin{equation}\label{eq:tdet}
\mu_{\mathrm{lot}}(T_{\mathrm{ref}})\,t_{\mathrm{det}}=\tau^*(h_0,N_0,N_{\mathrm{det}}),
\end{equation}
so the detection time combines the lot's own growth rate with its lag and starting load, which are precisely the quantities that dominate the error budget of Table~\ref{tab:anbudget}. Because the limit is reached at $\tau=K$, and $K-\tau^*(N_{\mathrm{det}})$ depends only weakly on $h_0$ and $N_0$ once the lag is over, a monitor that divides the logged integral by the measured detection time,
\begin{equation}\label{eq:cslcal}
S_{\mathrm{cal}}=\frac{1}{t_{\mathrm{det}}}\int_0^t\mu_{\mathrm{nom}}\big(T(s)\big)\dd s ,
\end{equation}
ranks trips by an estimate of $\tau/K$ in which the lot's rate scale and budget have been measured rather than assumed. What remains unmeasured is only the change of shape of $\mu(T)$ between $T_{\mathrm{ref}}$ and the storage temperature, which comes from the spread of $T_{\min}$. \secref{sec:calores} quantifies how much of the biological loss this recovers.


\begin{figure*}[!t]\centering\includegraphics{calorimetry}
\caption{Calorimetry and the generalized time. (a) Simulated heat-flow curves of the same milk biology (baseline parameters, $h_0=3.8$, $N_0=50$\,CFU\,mL$^{-1}$) under isothermal programmes at 20\degC{} (solid) and 12\degC{} (dashed) and a step from 6 to 22\degC{} after 30\,h (dotted), with the heat flow taken as proportional to the rate of cell formation. (b) The same curves divided by $\mu(T)$ and plotted against the generalized time $\tau$; they coincide (Proposition~\ref{prop:calo}). (c) AUC over the 12 baseline settings when the log is combined with information on the biology: none, the true lag and starting load, the true rate law, an isothermal calorimetric detection time at 20\degC{} with 5\,\% measurement error, and all of it (oracle). Bars span the settings and circles mark the mean}\label{fig:calorimetry}\end{figure*}

\subsection{Sampling error}
A sampled record replaces $\int\mu\dd t$ by a left Riemann sum. The error has a clean bound.

\begin{proposition}[Sampling bound]\label{prop:riemann}
Let readings be taken every $\Delta t$ and held until the next reading. (a) For any history, $|\CSL_{\Delta t}-\CSL|\le \Delta t\,V/K$, where $V$ is the total variation of $\mu(T(t))$ over the trip. (b) For a single rectangular excursion from $T_0$ to $T_1$, $|\CSL_{\Delta t}-\CSL|<\Delta t\,[\mu(T_1)-\mu(T_0)]/K$, whatever its timing. (c) If the excursion's timing is uniform relative to the sampling grid, the error has mean zero and standard deviation $\Delta t\,[\mu(T_1)-\mu(T_0)]/(K\sqrt6)$.
\end{proposition}

\begin{figure}[!htbp]\centering\includegraphics{sampling}
\caption{Why hourly readings are enough (Proposition~\ref{prop:riemann}). True growth rate during a 70-min warm spell to 13\degC{} (solid line) and the rate assumed from hourly readings held until the next reading (dashed line, dots). The hatched area on the rising side is missed and the dotted area on the falling side is over-counted; the two largely cancel}\label{fig:sampling}\end{figure}

Figure~\ref{fig:sampling} shows the mechanism on a single warm spell.

\begin{worked}
Hourly sampling of a 1-h excursion from 4 to 13\degC{} changes $\CSL$ by less than $(\AnMuThirteen-\AnMuFour)/\AnKnom=\AnRiem$, that is under 3\,\% of the shelf life, and on average by zero (standard deviation 0.011). Product thermal mass smooths the excursion further and makes the error smaller still.
\end{worked}


\subsection{Propagation of sensor error}
\begin{proposition}[Sensor propagation]\label{prop:sensor}
Let the sensor read $\tilde T=T+\beta+\varepsilon+e_q$ with a per-trip offset $\beta$, zero-mean noise $\varepsilon$ of variance $\sigma_n^2$ and quantisation error $e_q\sim\mathcal U(-\delta_q/2,\delta_q/2)$. To first order, for a trip near a steady temperature $\bar T$,
\begin{equation}\label{eq:sensor}
\Delta\ln\CSL\approx\frac{2\beta}{\bar T-T_{\min}},\qquad \mathrm{E}\,\Delta\CSL_{\varepsilon}=\frac{b^2}{K}\,\sigma_n^2\,t,
\end{equation}
and the noise and quantisation terms average out over $n$ readings as $1/\sqrt n$.
\end{proposition}
For a DS18B20 with offset uniform within $\pm0.5$\degC{} \cite{ds18b20}, $\mathrm{sd}(\beta)=0.5/\sqrt3=0.29$\degC, so $\mathrm{sd}(\ln\CSL)=\AnSdDS$; for a Class~A PT1000 within $\pm0.15$\degC{} \cite{iec60751}, \AnSdPT. Noise of 0.05\degC{} adds a bias of $\sigma_n^2/(\bar T-T_{\min})^2\approx2\times10^{-5}$, which is negligible.

\subsection{Propagation of biological variability}
Taking logarithms of $\CSL=(b^2/K)\int(T-T_{\min})^2\dd t$,
\begin{equation}\label{eq:lncsl}
\ln\CSL=2\ln b+\ln\!\int(T-T_{\min})^2\dd t-\ln K .
\end{equation}
Each unseen biological quantity enters additively on the log scale, and the delta method gives the contributions in Table~\ref{tab:anbudget}: a strain-to-strain $cv$ of 0.10 in $b$ gives $2\times0.10=\AnSdB$; a spread $\sigma_T=1.5$\,K in $T_{\min}$ gives $2\sigma_T/(\bar T-T_{\min})=\AnSdT$; and the spread of $h_0$ and $N_0$ gives \AnSdK{} through $K$. In quadrature the biology contributes \AnSdBio, about \AnRatio{} times the DS18B20 offset.

\begin{table}[!htbp]\centering\footnotesize
\caption{First-order error budget for $\ln\CSL$ at $\bar T=4$\degC{} (baseline biology).}\label{tab:anbudget}
\setlength{\tabcolsep}{4pt}
\begin{tabular}{@{}llc@{}}\toprule\rowcolor{dghead}
Source & Formula & sd of $\ln\CSL$\\\midrule
Growth-rate slope $b$ ($cv$ 0.10) & $2\,cv$ & \AnSdB\\
Minimum temperature ($\sigma_T$ 1.5\,K) & $2\sigma_T/(\bar T-T_{\min})$ & \AnSdT\\
Lag and starting load ($h_0$, $N_0$) & sd of $\ln K$ & \AnSdK\\
\textbf{Biology, combined} & quadrature & \textbf{\AnSdBio}\\\midrule
DS18B20 offset ($\pm0.5$\degC) & $2\,\mathrm{sd}(\beta)/(\bar T-T_{\min})$ & \AnSdDS\\
PT1000 offset ($\pm0.15$\degC) & same & \AnSdPT\\
Hourly sampling, 1-h excursion & Prop.~\ref{prop:riemann}(c) & 0.011\\
\bottomrule\end{tabular}\end{table}


\subsection{From estimation error to ranking accuracy}
Figure~\ref{fig:errors}a shows the budget of Table~\ref{tab:anbudget} graphically. Error in $\ln\CSL$ matters only through how often it swaps the order of a spoiled and an unspoiled trip. A simple model makes this precise.
\begin{proposition}[Binormal AUC]\label{prop:auc}
Let a trip's true log-margin be $z\sim\mathcal N(0,\sigma_z^2)$, spoiled iff $z>0$, and let the score be $s=z+e$ with independent $e\sim\mathcal N(0,\sigma_\mathrm{e}^2)$. Then, with $r=\sigma_e/\sigma_z$,
\begin{equation}\label{eq:binorm}
\AUC(r)=\frac12+\frac2\pi\arcsin\frac{1}{\sqrt{2(1+r^2)}} .
\end{equation}
\end{proposition}
$\AUC$ falls from 1 at $r=0$ to $\tfrac12$ as $r\to\infty$, and its slope flattens as $r$ grows. Inverting Eq.~\eqref{eq:binorm} at the simulated AUC of an ideal sensor gives a noise-to-signal ratio of about \AnR{} for biology alone. Adding the DS18B20 offset of Table~\ref{tab:anbudget} then moves $r$ only slightly, and Eq.~\eqref{eq:binorm} predicts a loss of \AnPredDS{} AUC; the simulation finds \AnObsDS. The model is cruder than the simulation, treating every error as Gaussian on one scale, but both put the sensor loss more than an order of magnitude below the biology loss.


\begin{figure*}[!t]\centering\includegraphics{errors}
\caption{Error budget and its effect on ranking. (a) First-order standard deviation of $\ln\CSL$ at 4\degC{} for each source (Table~\ref{tab:anbudget}); the upper four bars are biological, the lower three are instrumental. (b) Binormal relation between noise-to-signal ratio and AUC, Eq.~\eqref{eq:binorm}, with the oracle and a monitor with a perfect thermometer marked}\label{fig:errors}\end{figure*}

\subsection{Probabilistic output}
Rather than one verdict, the monitor can draw $M=80$ parameter sets $\theta_k=(b_k,T_{\min,k},h_{0,k},N_{0,k})$ from a prior and report
\begin{equation}\label{eq:mc}
\hat p=\frac1M\sum_{k=1}^M\mathbf 1\{\CSL(\theta_k)\ge1\},
\end{equation}
whose Monte Carlo error is $\mathrm{sd}(\hat p)=\sqrt{p(1-p)/M}\le0.056$.
A \emph{single alarm} is the special case of one draw at nominal values, $\hat p\in\{0,1\}$.


\section{Materials and methods}\label{sec:method}
\begin{figure*}[!t]\centering\includegraphics{design}
\caption{Design of the benchmark. Upper row: the simulated truth, hidden from the monitor, which turns a trip's temperature and biology into a true label. Lower row: what the monitor sees and computes from a sensed, sampled record. Removing information step by step gives the loss budget}\label{fig:design}\end{figure*}
\subsection{Overview}
Figure~\ref{fig:design} shows the design: a simulated world supplies the truth, and a monitor that sees only a sensed, sampled record tries to recover it. Figure~\ref{fig:flow} shows the pipeline from published data to the reported numbers. Every parameter has a recorded source (Table~\ref{tab:sources}); values without one are marked as assumed.

\begin{figure}[!htbp]\centering\includegraphics{flow}
\caption{Flowchart of the study, from trip generation and biology to sensing, scoring and statistical evaluation}\label{fig:flow}\end{figure}

\begin{table*}[!t]\centering\footnotesize
\caption{Where the data and parameters come from.}\label{tab:sources}
\begin{tabular}{@{}p{3.7cm}p{6.0cm}p{6.8cm}@{}}\toprule\rowcolor{dghead}
Source & What was taken & Used for\\\midrule
Neumeyer thesis \cite{neumeyer1995,neumeyer1997a} & $b=0.030$, $T_{\min}=-7.7$\degC{} from 68 \textit{Pseudomonas} growth curves in broth and milk & Baseline secondary model\\
Lin et al.\ \cite{lin2016} & Growth rates at 4, 15, 29\degC; lag at 4\degC; times to $10^6$\,CFU\,mL$^{-1}$ under two 4-h excursions & Cross-check of $b,T_{\min}$ ($b=0.027$, $T_{\min}=-9.0$\degC); $h_0$ range 2.5--5; validation\\
Zwietering et al.\ \cite{zwietering1994} & Extra lag after temperature shifts, about a quarter of the lag & Structure-risk run\\
Martin et al.\ \cite{martin2021}; Ranieri \& Boor \cite{ranieri2009} & Low post-pasteurisation contamination levels & Low-load run, $N_0$ 0.1--3\,CFU\,mL$^{-1}$\\
ComBase, spoilage flora \cite{combase} & \DatNSp{} maximum growth rates of aerobic spoilage bacteria in pasteurised milk, 2.7--14.9\degC & Milk-flora run; bias of isolate parameters\\
ComBase, \textit{L.~monocytogenes} \cite{combase} & \DatNLis{} records with a stated rate at 2--16\degC{} from 21 sources (\NumLisN{} fitting, \NumLisHeld{} held out) & \textit{Listeria} run; held-out validation\\
ComBase, milk compilation \cite{combase} & \DatNFull{} growth records in milk, nine organism groups & Open database (\secref{sec:data})\\
FSSAI Appendix B \cite{fssai2024} & Aerobic plate count $3\times10^4$\,CFU\,mL$^{-1}$ for pasteurised milk & Limit A; limit B $10^6$\,CFU\,mL$^{-1}$ follows \cite{lin2016}\\
DS18B20 datasheet \cite{ds18b20}; International Electrotechnical Commission (IEC) standard 60751 \cite{iec60751} & Accuracy $\pm0.5$\degC, step 0.0625\degC; PT1000 Class A $\pm0.15$\degC & Sensor models\\
Ndraha et al.\ \cite{ndraha2018}; G\"oransson et al.\ \cite{goransson2018} & Frequency and character of cold-chain breaks & Shape of the excursion model (rates assumed)\\
\bottomrule\end{tabular}\end{table*}

\subsection{Trip generator}
Each trip has a product-temperature trace on a one-minute grid over its horizon $H$. The ambient temperature is
\begin{equation}\label{eq:amb}
T_a(t)=T_{\mathrm{base}}+\sum_{k=1}^{N_e}(P_k-T_{\mathrm{base}})\,\mathbf 1\{s_k\le t<s_k+d_k\},
\end{equation}
with $T_{\mathrm{base}}\sim\mathcal U(2.5,5)$\degC; $N_e\sim\mathrm{Poisson}(H/144\,\mathrm h)$, capped at 8; start times $s_k$ uniform; peaks $P_k\in\{12,14.3,17.5,30.7\}$\degC{} with probabilities $\{0.30,0.25,0.30,0.15\}$; and durations $d_k=\mathrm{clip}(60\,\mathrm{e}^{0.8Z_k},10,480)$\,min with $Z_k\sim\mathcal N(0,1)$. The product follows with first-order thermal inertia,
\begin{equation}\label{eq:inertia}
\frac{\dd T}{\dd t}=\frac{T_a-T}{\tau_{\mathrm{th}}},
\end{equation}
with $\tau_{\mathrm{th}}\sim\mathcal U(5,20)$\,min for a low-mass pouch and $\mathcal U(1,4)$\,h for an insulated crate, plus a slow autoregressive wander $a_{j+1}=0.97a_j+\mathcal N(0,0.05^2)$, added as $0.1a_j$ (standard deviation about 0.02\degC).

\subsection{Reference biology}
A reference population grows along the trace by Eqs.~\eqref{eq:bar}--\eqref{eq:barq}, integrated by the Euler method at $\dd t=1$\,min; since $\mu\,\dd t\le0.03$ even at 30.7\degC, the integration error is negligible. Each trip draws $T_{\min}\sim\mathcal N(\bar T_{\min},\sigma_T^2)$, $b=\bar b\,\exp(cv\,Z-cv^2/2)$ with $cv\in\{0,0.10,0.25\}$, $h_0$ uniform over the measured range and $N_0$ uniform on the log scale (Table~\ref{tab:runs}). In the structure-risk run, whenever the product temperature departs from, or returns to, its two-hour running mean by more than 4\degC, the state is reset as $q\leftarrow q\,\mathrm{e}^{-0.25h_0}$, which re-imposes about a quarter of the lag \cite{zwietering1994}; the estimator is not told.

\begin{table}[!htbp]
\centering\footnotesize\caption{Benchmark runs. ``Assumed'' marks values without a source.}\label{tab:runs}
\begin{tabular}{@{}p{1.55cm}p{6.35cm}@{}}\toprule\rowcolor{dghead}
Run & Biology and labels\\\midrule
Baseline & $b=0.030$, $T_{\min}=-7.7$\degC{} \cite{neumeyer1995}; $\sigma_T=1.5$\,K (assumed); $h_0\in[2.5,5]$ \cite{lin2016}; $N_0$ 30--100\,CFU\,mL$^{-1}$ (assumed); limits $3\times10^4$, $10^6$; horizons \NumBaseHzA, \NumBaseHzB\,h\\
Extra lag & as baseline, plus extra lag after shifts \cite{zwietering1994}\\
Low load & as baseline, $N_0$ 0.1--3\,CFU\,mL$^{-1}$ \cite{martin2021}; horizons \NumLowHzA, \NumLowHzB\,h\\
Milk flora & $b=\NumCbB$, $T_{\min}=\NumCbT$\degC{} (ComBase, $n=\NumCbNfit$), low load; horizons \NumCombHzA, \NumCombHzB\,h\\
\textit{Listeria} & $b=\NumLisB$, $T_{\min}=\NumLisT$\degC{} (ComBase, $n=\NumLisN$); $h_0\in[2,4.5]$, $\sigma_T=1.0$\,K, $N_0$ 0.06--0.18\,CFU\,mL$^{-1}$ (assumed); labels $+0.5$, $+1$\,log; horizons \NumLisHzA, \NumLisHzB\,h\\\bottomrule
\end{tabular}
\end{table}

\subsection{Labels}
A trip is \emph{spoiled} if the reference population exceeds the limit within the horizon $H$. The horizon is the time the estimator takes to reach the limit at a constant 4\degC{} with nominal parameters, rounded up to a whole hour, as a printed shelf life would be; it was fixed before any method was compared. Prevalence is therefore constructed (about 0.5 for spoilage, \NumLisPrevMin--\NumLisPrevMax{} for \textit{Listeria}), and absolute predictive values must not be read as field rates. For \textit{Listeria} the labels are $+0.5$ and $+1$\,log$_{10}$ growth above the starting load; these are round-number labels, not regulatory limits.

\subsection{Sensors and sampling}
The sensed record is
\begin{equation}\label{eq:sensed}
\tilde T_i=\delta_q\,\mathrm{round}\!\left(\frac{T(t_i)+\beta+\varepsilon_i}{\delta_q}\right),
\end{equation}
with a per-trip offset $\beta\sim\mathcal U(-\beta_m,\beta_m)$, noise $\varepsilon_i\sim\mathcal N(0,\sigma_n^2)$ and $(\beta_m,\sigma_n,\delta_q)=(0.5,0.05,0.0625)$\degC{} for a DS18B20 \cite{ds18b20}, $(0.15,0.03,0.03125)$\degC{} for a PT1000 read by a MAX31865 converter \cite{iec60751}, and $(0,0,0)$ for an ideal sensor. Sampling intervals are 1, 15 and 60\,min plus an adaptive rule that relaxes to 15\,min while $|\Delta\tilde T|<0.2$\degC{} and $2\le\tilde T\le5$\degC, and returns to 1\,min otherwise.

\subsection{Scores}
Five scores are computed from each sensed record: $\CSL$ (Eq.~\eqref{eq:csl}); degree-hours above 4\degC, $\sum_i\Delta t_i(\tilde T_i-4)^+$; time above 5\degC; time above 8\degC; and the maximum temperature. An \emph{oracle} $\CSL$ uses each trip's true parameters and noise-free temperature and sets the ceiling. The probability of Eq.~\eqref{eq:mc} is computed under a \emph{matched} prior (the true spread), an \emph{overconfident} prior (no spread in $b$) and as a single alarm.

\subsection{Experimental design}
Each run has 12 settings (two limits $\times$ two thermal regimes $\times$ three values of $cv$) with 4\,000 trips each, so five runs comprise 240\,000 trips. Every trip is read by three sensors at four sampling rules, and every record is scored five ways, so all comparisons are \emph{paired}: they differ only in the factor of interest. The loss budget removes information in a fixed order (oracle $\to$ nominal biology with ideal 1-min record $\to$ PT1000 at 15\,min $\to$ DS18B20 at 15 and 60\,min). Random seeds are fixed; Table~\ref{tab:design} summarises the design.

\begin{table}[!htbp]\centering\footnotesize
\caption{Experimental design.}\label{tab:design}
\begin{tabular}{@{}ll@{}}\toprule\rowcolor{dghead}
Factor & Levels\\\midrule
Parameter set (run) & 5 (Table~\ref{tab:runs})\\
Limit & 2 ($3\times10^4$, $10^6$\,CFU\,mL$^{-1}$; or $+0.5$, $+1$\,log)\\
Thermal regime & 2 (pouch, crate)\\
Strain spread $cv$ & 3 (0, 0.10, 0.25)\\
Trips per setting & 4\,000 (total 240\,000)\\
Sensor & 3 (ideal, PT1000, DS18B20)\\
Sampling & 4 (1, 15, 60\,min, adaptive)\\
Score & 5 + oracle + 3 probability variants\\
Bootstrap & $B=200$ (single), $B=400$ (paired)\\
\bottomrule\end{tabular}\end{table}

\subsection{Benchmark protocol}\label{sec:protocol}
The benchmark is a factorial simulation experiment. For every run and setting it executes the procedure below; every random number is drawn from a generator with a fixed seed, so the whole study is reproducible to the last digit.

\begin{tcolorbox}[enhanced,breakable,colback=white,colframe=dgrule,boxrule=0.5pt,arc=0pt,left=5pt,right=5pt,top=3pt,bottom=3pt,fontupper=\small,title={\small\bfseries Algorithm 1\enspace Benchmark for one run and one setting},coltitle=black,colbacktitle=dgpale]
\begin{enumerate}[leftmargin=1.6em,label=\arabic*.]
\item Set the limit $N_L$ and the horizon $H=\lceil\SL_{\mathrm{nom}}(4\,^\circ\mathrm C)\rceil$.
\item For $i=1,\dots,4000$: draw $T_{\mathrm{base}}$, the excursions $(s_k,d_k,P_k)$ and $\tau_{\mathrm{th}}$; build $T_a(t)$ by Eq.~\eqref{eq:amb} and the product trace $T_i(t)$ by Eq.~\eqref{eq:inertia} on a 1-min grid.
\item Draw the trip's biology $(b_i,T_{\min,i},h_{0,i},N_{0,i})$; integrate Eqs.~\eqref{eq:bar}--\eqref{eq:barq} along $T_i(t)$; set $y_i=1$ if $N_i(H)>N_L$.
\item For each sensor $\in\{$ideal, PT1000, DS18B20$\}$ and interval $\in\{1,15,60$\,min, adaptive$\}$: form the sensed record by Eq.~\eqref{eq:sensed}.
\item On each record compute CSL, degree-hours, time above 5 and 8\degC, maximum temperature, the Arrhenius and $Q_{10}$ integrators of Eq.~\eqref{eq:mkt}, and the Monte Carlo probability of Eq.~\eqref{eq:mc}; compute the oracle CSL from the true trace and biology.
\item Compute every AUC with its bootstrap interval, every paired contrast with its paired-bootstrap interval, and the Brier score and ECE of the probabilities (\secref{sec:stats}).
\end{enumerate}
\end{tcolorbox}

\subsection{Parameter estimation from ComBase}
ComBase reports maximum rates in log$_{10}$ units; we convert with $\mu=\mu_{10}\ln10$ and fit Eq.~\eqref{eq:ols}. For \textit{Listeria}, records with no added factors form the fitting set and records with other stated conditions (sterilised or previously heated milk, added salts or acids) form the held-out set; we report $B_f$ and $A_f$ (Eq.~\eqref{eq:bfaf}) on both.

\section{Statistical methodology}\label{sec:stats}
The study makes three kinds of statistical statement: how well a score ranks spoiled against unspoiled trips, whether one score ranks better than another on the same trips, and how well a probability is calibrated. A fourth kind, how well a temperature law fits measured growth rates, was used above. Figure~\ref{fig:statflow} summarises the procedures; this section defines them and justifies the sample sizes.

\begin{figure}[!htbp]\centering\includegraphics{statflow}
\caption{Statistical procedures of the study, from the per-trip labels and scores to point estimates, intervals, paired comparisons, calibration and growth-rate fits}\label{fig:statflow}\end{figure}

\subsection{The estimand: area under the ROC curve}
Let $S_1$ be the score of a randomly chosen spoiled trip and $S_0$ that of a randomly chosen unspoiled trip. The quantity of interest is
\begin{equation}\label{eq:aucdef}
\AUC=\Pr(S_1>S_0)+\tfrac12\Pr(S_1=S_0),
\end{equation}
the probability that the monitor ranks a spoiled trip above an unspoiled one \cite{hanley1982}. It does not depend on a threshold or on the prevalence, which in this benchmark is constructed, so it is the natural measure of how much information a score carries. For $n_1$ spoiled and $n_0$ unspoiled trips with scores $s$, its unbiased estimate is the Mann--Whitney statistic divided by $n_1n_0$ \cite{mannwhitney1947},
\begin{equation}\label{eq:auc}
\widehat{\AUC}=\frac{1}{n_1n_0}\sum_{i:y_i=1}\sum_{j:y_j=0}\Big[\mathbf 1\{s_i>s_j\}+\tfrac12\mathbf 1\{s_i=s_j\}\Big],
\end{equation}
computed in $O(n\log n)$ from the ranks $r_i$ of the pooled scores as $\widehat{\AUC}=(\sum_{y_i=1}r_i-n_1(n_1+1)/2)/(n_1n_0)$.

\subsection{Precision and sample size}
Hanley and McNeil approximated the standard error of $\widehat{\AUC}$ by
\begin{equation}\label{eq:hm}
\mathrm{SE}^2=\frac{A(1-A)+(n_1-1)(Q_1-A^2)+(n_0-1)(Q_2-A^2)}{n_1n_0},
\end{equation}
with $Q_1=A/(2-A)$ and $Q_2=2A^2/(1+A)$ \cite{hanley1982}. For $A=0.67$ and $n_1=n_0=2000$, the case of a typical setting with 4\,000 trips and prevalence near one half, $\mathrm{SE}=0.0085$, so a single AUC is known to about $\pm0.017$. This fixed the number of trips: the biology gap of 0.27--0.40 is more than fifteen standard errors, while the sensor and sampling effects of 0.01--0.02 are of the order of one standard error and can be resolved only by paired comparison.

\subsection{Interval estimates by the bootstrap}
Intervals for single AUCs are percentile intervals from $B=200$ bootstrap resamples of trips \cite{efron1993}: trips are drawn with replacement, Eq.~\eqref{eq:auc} is recomputed, and the 2.5th and 97.5th percentiles of the replicates are reported. The trip is the resampling unit because trips are independent by construction, while the many readings within a trip are not.

\subsection{Paired comparisons}
Two scores computed on the same trips are correlated, because both respond to the same temperature histories. For a difference $D=\AUC_1-\AUC_2$,
\begin{equation}\label{eq:vard}
\mathrm{Var}(\hat D)=\mathrm{Var}(\widehat{\AUC}_1)+\mathrm{Var}(\widehat{\AUC}_2)-2\,\mathrm{Cov}(\widehat{\AUC}_1,\widehat{\AUC}_2),
\end{equation}
and the covariance term is large and positive for similar scores. Testing the two AUCs separately would ignore it and lose most of the power; this is the reason the benchmark reads every trip with every sensor, sampling rule and score (a common-random-numbers design). The interval for $D$ is obtained by resampling the same trips for both scores, $B=400$ times, and taking the percentiles of the replicate differences. This paired bootstrap is the resampling counterpart of the method of DeLong et al.\ \cite{delong1988}. Figure~\ref{fig:stats}c shows the bootstrap distribution for the kinetic score against degree-hours in one setting: the difference is only 0.019, yet its 95\,\% interval excludes zero because the two scores are strongly correlated.

\subsection{Multiplicity and aggregation over settings}
Each run has 12 settings and each contrast is examined in all of them. We do not draw a conclusion from any single setting. Instead we report, for each contrast, the range of the difference and the number of settings whose interval excludes zero. If a contrast had no real effect, each interval would exclude zero in the favourable direction with probability 0.025, so among 60 settings about 1.5 such exclusions would be expected by chance; observing many more is evidence of a real effect. As a further check, a sign test counts the settings in which one score is ahead. The settings share parameter values and are therefore not fully independent, so the binomial probabilities quoted with the sign test are indicative rather than exact.

\subsection{Calibration of probabilities}
A probability $\hat p_i$ is calibrated if, among trips given probability $p$, a fraction $p$ is spoiled. We report the Brier score and the expected calibration error over ten equal-width bins $B_m$ \cite{brier1950,naeini2015},
\begin{gather}
\mathrm{Brier}=\frac1n\sum_i(\hat p_i-y_i)^2,\label{eq:brier}\\
\mathrm{ECE}=\sum_{m=1}^{10}\frac{|B_m|}{n}\,\big|\bar y_{B_m}-\bar p_{B_m}\big|,\label{eq:ece}
\end{gather}
and show reliability diagrams (Fig.~\ref{fig:stats}b), which plot $\bar y_{B_m}$ against $\bar p_{B_m}$. The Monte Carlo probability is itself an estimate: with $M=80$ parameter draws its standard error is $\sqrt{p(1-p)/M}\le0.056$, small compared with the miscalibration of a single alarm.

\subsection{Growth-rate fits}
Secondary models were fitted by ordinary least squares, Eq.~\eqref{eq:ols}, with 95\,\% intervals for $b$ and $T_{\min}$ from 2\,000 bootstrap resamples of records. Fits are compared by the small-sample Akaike criterion
\begin{equation}\label{eq:aicc}
\mathrm{AICc}=n\ln\frac{\mathrm{RSS}}{n}+2k+\frac{2k(k+1)}{n-k-1},
\end{equation}
where RSS is the residual sum of squares of $\ln\mu$ and $k$ the number of parameters \cite{akaike1974,hurvich1989}; by cross-validation in which all records at one temperature are left out in turn; and by the bias and accuracy factors of Ross \cite{ross1996},
\begin{equation}\label{eq:bfaf}
B_f=10^{\frac1n\sum\log_{10}(\hat\mu_i/\mu_i)},\qquad A_f=10^{\frac1n\sum|\log_{10}(\hat\mu_i/\mu_i)|},
\end{equation}
where $B_f=1$ means no bias and $A_f-1$ is the typical relative error.

\begin{figure*}[!t]\centering\includegraphics{stats}
\caption{Statistics of one representative setting (baseline biology, $3\times10^4$\,CFU\,mL$^{-1}$, low-mass pouch, $cv=0.10$, DS18B20 sampled every 15\,min, 4\,000 trips). (a) ROC curves of the oracle (dotted), the kinetic score CSL (solid), degree-hours (dashed), time above 5\degC{} (fine dotted) and maximum temperature (dash-dotted), with their AUC. (b) Reliability diagram of the Monte Carlo probability (circles) and the single alarm (squares) against perfect calibration (dashed). (c) Bootstrap distribution of the paired AUC difference between CSL and degree-hours, with its 95\,\% interval (dashed) and zero (solid)}\label{fig:stats}\end{figure*}



\section{The open data}\label{sec:data}
We compiled every ComBase growth record whose food matrix is milk, \DatNFull{} in all, across nine organism groups (Table~\ref{tab:dcomp}). For each record we keep the organism, food description, temperature, pH, water activity, other stated conditions, the maximum growth rate where the record states one, and the laboratory or publication it comes from. A rate is stated for 179 records; 131 give only a growth curve that would need fitting, 20 report no growth, and for 43 the rate was not transcribed. All records are kept and flagged by status, since curve and no-growth records are useful for growth/no-growth modelling even though we do not use them here. Two sets drive this study: \DatNSp{} records of aerobic spoilage flora in pasteurised milk from one laboratory, retrieved by a separate search of the same database (listed in full in the appendix), and the \DatNLis{} records of \textit{L.~monocytogenes} at 2--16\degC{} from 21 sources (\ListeriaRef). All are released in the repository as comma-separated (CSV) files, so that every number in \secref{sec:val} can be traced to the record it came from.

\TabDataComp

\section{Validation against real data}\label{sec:val}
Three checks against real records are reported below and summarised in Table~\ref{tab:valsum}.
\subsection{A published fluctuating-temperature experiment}
Lin et al.\ held UHT milk inoculated with \textit{P.~fluorescens} at 4\degC{} with two four-hour excursions to 29\degC, and report that time to $10^6$\,CFU\,mL$^{-1}$ fell from about 90\,h to about 55\,h \cite{lin2016}. With parameters fitted only to their growth rates and their 4\degC{} lag, our reference model gives \NumLinConst\,h and \NumLinExc\,h, an error of $\NumLinErr$\,\% on the excursion case, which was not used for fitting. The published times are approximate values from the text; this is one experiment with one organism and an inoculated load far above real post-pasteurisation levels.

\subsection{Natural spoilage flora}
The \NumCbN{} ComBase records of aerobic spoilage bacteria in pasteurised milk span 2.7--14.9\degC. One record (6.8\degC) lies far below its neighbours; without it the square-root law fits the remaining \NumCbNfit{} records with $b=\NumCbB$, $T_{\min}=\NumCbT$\degC{} and a residual standard deviation of \NumCbSd{} (natural-log units, including experimental noise), and with it $b=\NumCbAllB$ and $T_{\min}=\NumCbAllT$\degC, so the conclusions do not depend on it. Parameters from inoculated isolates overpredict these records with a bias factor of \NumCbBias{} that is nearly constant across temperature (Fig.~\ref{fig:cb}): the temperature \emph{shape} transfers, the \emph{level} does not. This is consistent with studies that choose a fast isolate as a worst case.

\begin{figure}[!htbp]\centering\includegraphics{combase}
\caption{Maximum growth rates of natural spoilage flora in pasteurised milk (points; data from ComBase \cite{combase}) with the Neumeyer-anchored model (dotted line), the Lin-anchored model (dashed line) and the square-root fit to the records (solid line)}\label{fig:cb}\end{figure}

\subsection{\textit{Listeria monocytogenes}}
The fit to the \NumLisN{} fitting records gives $b=\NumLisB$ (95\,\% bootstrap interval \NumLisBLo--\NumLisBHi) and $T_{\min}=\NumLisT$\degC{} ($\NumLisTLo$ to $\NumLisTHi$\degC), a generation time of \NumLisGen\,h at 4\degC, and $B_f=\NumLisBiasHeld$, $A_f=\NumLisAccHeld$ on the \NumLisHeld{} held-out records ($A_f=\NumLisAccFit$ on the fitted ones). The split is by condition, not by study, so some records share a source with the fitted ones and the test is weaker than a split by study. Because $T_{\min}$ is closer to refrigeration temperature, Eq.~\eqref{eq:q10} makes the rate at 13\degC{} \NumLisRatio{} times the rate at 4\degC{} for \textit{Listeria}, against about 3.4 times for spoilage flora (Fig.~\ref{fig:lis}).

\begin{figure*}[!t]\centering\includegraphics{lis}
\caption{\textit{Listeria monocytogenes} in milk. (a) Maximum growth rates used for fitting (filled circles) and held out by condition (open circles), with the square-root fit (data from ComBase \cite{combase}). (b) Growth rate relative to 4\degC{} for spoilage flora (solid line) and \textit{Listeria} (dashed line)}\label{fig:lis}\end{figure*}

\begin{table}[!htbp]\centering\footnotesize
\caption{Validation summary.}\label{tab:valsum}
\setlength{\tabcolsep}{3pt}
\begin{tabular}{@{}lll@{}}\toprule\rowcolor{dghead}
Check & Published or observed & This model\\\midrule
Lin, constant 4\degC & about 90\,h & \NumLinConst\,h\\
Lin, two excursions & about 55\,h & \NumLinExc\,h ($\NumLinErr$\,\%)\\
Spoilage flora, $n=\NumCbNfit$ & $b=\NumCbB$, $T_{\min}=\NumCbT$ & isolates: $B_f=\NumCbBias$\\
\textit{Listeria}, held out, $n=\NumLisHeld$ & -- & $B_f=\NumLisBiasHeld$, $A_f=\NumLisAccHeld$\\
Generation time at 4\degC & \textit{Pseudomonas} 5.6\,h \cite{neumeyer1995} & 5.6\,h; \textit{Listeria} \NumLisGen\,h\\
\bottomrule\end{tabular}\end{table}

\section{Results}\label{sec:results}
\begin{keyresults}
\begin{itemize}[leftmargin=1.1em]
\item \textbf{Biology dominates.} Unseen biology costs \NumBaseGapMin--\NumBaseGapMax{} AUC in the baseline run (\NumLisGapMin--\NumLisGapMax{} for \textit{Listeria}); no sensor or sampling rate recovers it.
\item \textbf{Instruments matter little.} A PT1000 instead of a DS18B20 gains at most \NumBaseSensMax{} AUC; hourly sampling is indistinguishable from minute sampling in \NumBaseSampNull{} of 12 baseline settings.
\item \textbf{The square-root integrator is never outperformed.} Against degree-hours, Arrhenius (mean kinetic temperature) and $Q_{10}$ integrators it is ahead in most settings and behind in none, by about 0.01 AUC.
\item \textbf{Calorimetry recovers most of the biology.} An isothermal calorimetric detection time of the lot, combined with the log, closes \CalFracMin--\CalFracMax\,\% of the gap to the oracle (AUC \CalCalMin--\CalCalMax{} against \CalNomMin--\CalNomMax).
\item \textbf{Probabilities beat alarms.} A Monte Carlo probability has ECE \NumBaseEceMMin--\NumBaseEceMMax{} against \NumBaseEcePMin--\NumBaseEcePMax{} for a single alarm.
\end{itemize}
\end{keyresults}
\subsection{Loss budget (RQ1)}
An oracle that knows each trip's true biology and temperature separates spoiled from unspoiled trips almost perfectly, as Proposition~\ref{prop:csl} predicts. The same $\CSL$ with nominal biology and an ideal sensor loses \NumBaseGapMin--\NumBaseGapMax{} AUC across the 12 settings of the baseline run (Table~\ref{tab:budget}, Fig.~\ref{fig:budget}). The loss is large even when strains do not differ in growth rate ($cv=0$), because $T_{\min}$, lag state and starting load still differ between trips. None of it is in the temperature record, so no sensor or sampling rate can recover it.

\TabBudget

\begin{figure}[!htbp]\centering\includegraphics{budget}
\caption{Loss budget for (a) a low-mass pouch and (b) an insulated crate (limit $3\times10^4$\,CFU\,mL$^{-1}$, $cv=0.10$). The top row is the oracle with true biology and true temperature; the rows below remove information step by step. Bars are 95\,\% bootstrap intervals}\label{fig:budget}\end{figure}

\subsection{Against simple rules (RQ3)}
Table~\ref{tab:auc} in the Appendix gives the AUC of every score in every baseline setting. On identical trips the kinetic score beats time above 5\degC{} in \NumBaseFiveSig{} of 12 settings and maximum temperature in \NumBaseMaxSig. Against degree-hours above 4\degC, a rule that needs no growth model, the gain is positive with an interval excluding zero in \NumBaseGainSig{} of 12 settings but small: \NumBaseGainMin--\NumBaseGainMax{} AUC in the baseline run and up to \NumLagGainMax{} when extra lag is present (Table~\ref{tab:contrasts}). Corollary~\ref{cor:dh} explains why: degree-hours are the first-order term of the same clock.


\subsection{Sensors and sampling (RQ2)}
Replacing a DS18B20 with a PT1000 gains \NumBaseSensMin--\NumBaseSensMax{} AUC at 15-minute sampling. Sampling hourly instead of every minute shows no detectable difference in \NumBaseSampNull{} of 12 settings; where it does, the effect is at most about 0.01. The adaptive rule is indistinguishable from fixed one-minute sampling. These results depend on the assumed excursion model (minimum duration 10\,min, median 60\,min) and would change if short, sharp excursions were common, as Proposition~\ref{prop:riemann} shows.

\TabContrasts

\subsection{Robustness across parameter sets}
Figure~\ref{fig:runs} summarises the four effects across all five runs. The pattern is the same whether the biology is anchored to inoculated isolates, includes extra lag, starts from a realistic low load, is fitted to natural flora, or describes \textit{Listeria}: the biology gap is \NumLowGapMin--\NumLowGapMax{} (low load), \NumCombGapMin--\NumCombGapMax{} (milk flora) and \NumLisGapMin--\NumLisGapMax{} (\textit{Listeria}); the gain over degree-hours stays below \NumLagGainMax{} in every run, and the sensor gain stays small (\NumLisSensMin--\NumLisSensMax{} for \textit{Listeria}).

\begin{figure*}[!t]\centering\includegraphics{runs}
\caption{Four effects across the five parameter sets, as AUC differences on identical trips. (a) Loss to unobserved biology (oracle minus best monitor). (b) Gain of the kinetic score over degree-hours. (c) Gain of a PT1000 over a DS18B20. (d) Change from hourly instead of every-minute sampling. Positive values favour the first option. Bars span the 12 settings of each run and open circles mark the median; horizontal scales differ}\label{fig:runs}\end{figure*}

\subsection{Against Arrhenius and $Q_{10}$ integrators (RQ3)}\label{sec:rivalres}
The rules most often used in practice to summarise a temperature log are a mean kinetic temperature, that is an Arrhenius integrator with a fixed activation energy, and a $Q_{10}$ rule. We compared the square-root integrator CSL with three of them on identical trips: Arrhenius with the activation energy fitted to the ComBase records of the run's organism (\SecEaSp\,kJ\,mol$^{-1}$ for spoilage flora, \SecEaLis\,kJ\,mol$^{-1}$ for \textit{Listeria}), Arrhenius with the 120\,kJ\,mol$^{-1}$ measured calorimetrically for a food-spoilage process \cite{wadso2011}, and $Q_{10}=3$ \cite{zou2023} (Table~\ref{tab:rivals}, Fig.~\ref{fig:rivals}). No integrator was significantly ahead of CSL in any of the 60 settings. CSL was ahead of the fitted Arrhenius integrator in \RvFitAhead{} of 60 settings, significantly so in \RvFitSig; ahead of the 120\,kJ\,mol$^{-1}$ integrator in \RvWadAhead{} (\RvWadSig{} significant); and ahead of the $Q_{10}$ rule in \RvQtenAhead{} (\RvQtenSig{} significant). By the sign test these proportions are far from what chance would give ($p<10^{-5}$ for each), with the caveat on dependence noted in \secref{sec:stats}. The advantage is largest when extra lag follows temperature shifts, where CSL leads in every setting by up to 0.038 AUC, and it grows with the activation energy, as Table~\ref{tab:weights} predicts. The margins are nevertheless small, of the order of 0.01 AUC: every integrator is a generalized time, and all are limited far more by unobserved biology than by their temperature law.

\TabRivals

\begin{figure*}[!t]\centering\includegraphics{rivals}
\caption{Paired AUC differences between the square-root integrator CSL and (a) an Arrhenius integrator with the activation energy fitted to ComBase records, (b) an Arrhenius integrator with $E_a=120$\,kJ\,mol$^{-1}$, (c) a $Q_{10}=3$ rule and (d) degree-hours, on identical trips (DS18B20, 15\,min). Bars span the 12 settings of each run and open circles mark the median; the numbers give the settings whose 95\,\% paired-bootstrap interval excludes zero. Positive values favour CSL}\label{fig:rivals}\end{figure*}

\subsection{Calibration (RQ4)}
A single alarm using nominal parameters is badly calibrated (ECE \NumBaseEcePMin--\NumBaseEcePMax{}, Table~\ref{tab:cal}). A Monte Carlo probability with a matched prior reaches \NumBaseEceMMin--\NumBaseEceMMax, and an overconfident prior \NumBaseEceOMin--\NumBaseEceOMax. Absolute values reflect the constructed prevalence; the ranking of the three variants is the finding. The \textit{Listeria} run needs a caveat. With the $+1$\,log label the probability reaches ECE \NumLisEceBMMin--\NumLisEceBMMax{} against \NumLisEceBPMin--\NumLisEceBPMax{} for the alarm. With the half-log label the limit lies inside the early lag-to-growth transition, where the Gompertz estimator is a poor approximation of the reference; even the oracle reaches only \NumLisOrcAMin--\NumLisOrcAMax{} AUC there (\NumLisOrcBMax{} for the $+1$\,log label), and the probability, though still better than the alarm (ECE \NumLisEceAMMin--\NumLisEceAMMax{} against \NumLisEceAPMin--\NumLisEceAPMax), is poorly calibrated in absolute terms. Label choice interacts with estimator structure.

\TabCal

\subsection{Value of a calorimetric measurement of the lot (RQ5)}\label{sec:calores}
The loss budget says what is lost; Proposition~\ref{prop:calo} suggests how much of it a laboratory measurement could buy back. We regenerated the 12 baseline settings with a separate random stream (4\,000 trips each) and scored every trip from the same DS18B20 record, sampled every 15\,min, with four kinds of added information (Table~\ref{tab:calo}, Fig.~\ref{fig:calorimetry}c). Knowing each trip's true lag and starting load raises the AUC by only \CalGainBudMin--\CalGainBudMax. Knowing its true rate law ($b$ and $T_{\min}$) raises it by \CalGainRateMin--\CalGainRateMax, to \CalRateMin--\CalRateMax. This ordering is the one predicted by the first-order budget of Table~\ref{tab:anbudget}, in which the rate parameters contribute \AnSdB{} and \AnSdT{} to the spread of $\ln\CSL$ and lag and load only \AnSdK.

A calorimetric detection time combines rate, lag and load in one number, Eq.~\eqref{eq:tdet}. We simulated an isothermal measurement of each trip's biology at $T_{\mathrm{ref}}=20$\degC{} to a detection level of $10^5$\,CFU\,mL$^{-1}$, with a lognormal measurement error of 5\,\%, and scored the trip by Eq.~\eqref{eq:cslcal}. The detection times have a median of \CalTdetMed\,h (5th--95th percentile \CalTdetLo--\CalTdetHi\,h), short enough for a sample taken at dispatch. The score reaches AUC \CalCalMin--\CalCalMax, against \CalNomMin--\CalNomMax{} for the log alone; the gain exceeds zero in all \CalGainCalSig{} settings, and the measurement closes \CalFracMin--\CalFracMax\,\% (mean \CalFracMean\,\%) of the gap to the oracle. With a 15\,\% measurement error the AUC is \CalSensBMin--\CalSensBMax{} (mean \CalSensBFrac\,\% of the gap), and with $T_{\mathrm{ref}}=10$\degC{} it is \CalSensCMin--\CalSensCMax{} (mean \CalSensCFrac\,\%). What remains, \CalRestMin--\CalRestMax{} AUC, is the change of shape of $\mu(T)$ between the reference and storage temperatures, which a single isothermal measurement cannot resolve.

These numbers are an upper bound in one respect: each simulated trip carries the biology of its own sample, whereas in practice a sample represents a lot whose cartons differ among themselves. They also assume a constant heat yield per cell and a fixed detection count. Within these assumptions the result is clear: one isothermal calorimetric measurement is worth an order of magnitude more ranking accuracy than any improvement of the thermometer.

\TabCalo

\subsection{Mathematics and simulation agree}
Table~\ref{tab:agree} sets each analytic prediction of \secref{sec:math} beside the corresponding simulated result. They were obtained independently: the propositions use only the model equations and parameter ranges, and the simulation uses none of the propositions.

\begin{table*}[!t]\centering\footnotesize
\caption{Analytic predictions (\secref{sec:math}) against simulated results (baseline run unless stated).}\label{tab:agree}
\begin{tabular}{@{}p{4.3cm}p{5.6cm}p{6.6cm}@{}}\toprule\rowcolor{dghead}
Question & Mathematics predicts & Simulation finds\\\midrule
Does an oracle rank perfectly? & Yes, both models read $\tau$ (Props.~\ref{prop:clock}, \ref{prop:csl}) & Oracle AUC 1.000 (baseline); below 1 only where structure differs (extra lag 0.988, \textit{Listeria} half-log)\\
Biology versus sensor & sd of $\ln\CSL$: \AnSdBio{} against \AnSdDS, ratio about \AnRatio & Biology gap \NumBaseGapMin--\NumBaseGapMax{} AUC against DS18B20 loss of about 0.01\\
Effect of sensor on AUC & Binormal model: \AnPredDS{} (Prop.~\ref{prop:auc}) & \AnObsDS{} (ideal $\to$ DS18B20, 1\,min)\\
Degree-hours versus CSL & First-order term of the same generalized time (Cor.~\ref{cor:dh}) & CSL ahead by at most \NumLagGainMax{} AUC\\
Hourly versus every-minute & Error $<\AnRiem$ per 1-h excursion, mean zero (Prop.~\ref{prop:riemann}) & $\NumBaseSampMin$ to $\NumBaseSampMax$ AUC; not detectable in \NumBaseSampNull{} of 12\\
Which biology matters most & Rate parameters: sd \AnSdB{} ($b$) and \AnSdT{} ($T_{\min}$) against \AnSdK{} (lag, load) & True rate law adds \CalGainRateMin--\CalGainRateMax{} AUC; true lag and load \CalGainBudMin--\CalGainBudMax\\
Calorimetric measurement & Detection time measures rate and budget together (Prop.~\ref{prop:calo}) & Closes \CalFracMin--\CalFracMax\,\% of the biology gap\\
Value of a probability & sd of $\hat p\le0.056$ with $M=80$ & ECE \NumBaseEceMMin--\NumBaseEceMMax{} against \NumBaseEcePMin--\NumBaseEcePMax{} for an alarm\\
\bottomrule\end{tabular}\end{table*}

\section{DairyGuard toolkit and demonstrator}\label{sec:toolkit}
The study is delivered as a working toolkit rather than a set of scripts.

\subsection{Monitor architecture}
Figure~\ref{fig:system} shows the architecture we propose for a logger. A temperature probe feeds a sampling scheduler that applies the adaptive rule; each accepted reading is converted to a shelf-life fraction by a look-up table of $1/\SL(T)$ and added to the accumulator (Eq.~\eqref{eq:csl}); a probability module evaluates Eq.~\eqref{eq:mc} over stored parameter draws; and the result is shown as a risk band and sent over Bluetooth Low Energy. Parameters are estimated offline from data such as ComBase and loaded into the table. The state that must persist between readings is a single number per parameter draw, which suits a microcontroller.

\begin{figure*}[!t]\centering\includegraphics{system}
\caption{Block diagram of a DairyGuard monitor, read in the order of the arrows (steps 1--6), with the value each step would hold 40\,min after a truck door was left open. Steps 2--5 run on the logger; the shelf-life table and the parameter sets are prepared once, offline}\label{fig:system}\end{figure*}

\subsection{Software}
The repository is organised in layers. A kinetics module implements Eqs.~\eqref{eq:sqrt}--\eqref{eq:csl}; a simulator implements \secref{sec:method}; the benchmark computes scores, AUCs, paired bootstraps and calibration; validation scripts reproduce \secref{sec:val}; and every number in this paper is written into the manuscript by a script from the result files. Numerical work uses NumPy and SciPy \cite{harris2020,virtanen2020}, and all graphs and diagrams were drawn with Matplotlib \cite{hunter2007}. Fourteen automated tests, listed in Table~\ref{tab:tests} of the Appendix, run on every change. Parts of the simulation and analysis code were written with the help of an AI coding assistant; every function is covered by the automated tests and was checked by the authors.

\subsection{Browser demonstrator}
A public web demonstrator runs the same kernel in the browser (Fig.~\ref{fig:demo}); a parity test guarantees that it gives the same numbers as the Python code. A user picks a trip, such as a shop door left open or a failed truck refrigerator, or draws one, and sees the probability that the milk has passed the limit, the 60 cartons of a truck coloured by their simulated fate, and the comparison with a simple alarm. The demonstrator is software only; a hardware logger is future work.

\begin{figure}[!htbp]\centering\includegraphics[width=84mm]{demo.jpg}
\caption{The browser demonstrator for a trip in which the truck refrigerator failed: the monitor reports a probability of passing the limit instead of a single alarm}\label{fig:demo}\end{figure}

\section{Discussion}\label{sec:disc}
\subsection{Implications for monitor design}
In every run, what the thermometer cannot see costs about an order of magnitude more accuracy than any choice of sensor or sampling interval, and \secref{sec:math} shows that this is a structural property rather than an accident of the simulation: the log is reduced to one number, biology enters that number additively on the log scale, and its spread is about seven times that of a cheap probe. A better thermometer is therefore a poor way to improve a spoilage estimate; a cheap digital probe sampled every 15--60\,min loses little under our assumptions. Effort is better spent on biology: measuring starting load or a handling-history proxy, learning the lag and strain spread for the product at hand, and reporting uncertainty instead of a single verdict. Table~\ref{tab:rules} turns this into design rules.

\begin{table}[!htbp]\centering\footnotesize
\caption{Design rules for a milk cold-chain monitor.}\label{tab:rules}
\begin{tabular}{@{}p{2.3cm}p{5.6cm}@{}}\toprule\rowcolor{dghead}
Decision & Rule and evidence\\\midrule
Sensor & A $\pm0.5$\degC{} digital probe suffices; a PT1000 gains $\le$\NumBaseSensMax{} AUC (Prop.~\ref{prop:sensor})\\
Sampling & 15--60\,min is enough unless excursions under 10\,min are common (Prop.~\ref{prop:riemann})\\
Score & $\CSL$; degree-hours is an acceptable fallback (Cor.~\ref{cor:dh})\\
Output & A probability, not an alarm (ECE \NumBaseEceMMin--\NumBaseEceMMax{} vs \NumBaseEcePMin--\NumBaseEcePMax)\\
Parameters & Fit to natural flora; isolates overpredict by \NumCbBias\\
Improvement & Measure biology, above all the rate law of the lot\\
Laboratory test & Isothermal calorimetric detection time at dispatch; closes \CalFracMin--\CalFracMax\,\% of the biology gap (Prop.~\ref{prop:calo})\\
\bottomrule\end{tabular}\end{table}

\subsection{A role for calorimetry in the cold chain}
Calorimetry cannot travel with a crate, and a temperature log cannot see the biology; the two are complementary. Proposition~\ref{prop:calo} shows that they measure the same generalized time from opposite sides: the log measures how much of it has elapsed, and an isothermal detection time measures, in the same units, how much of it the lot can afford. In the benchmark a single detection time at dispatch recovers most of the loss that no sensor can touch (\secref{sec:calores}). Isothermal microcalorimeters already detect contamination in hours \cite{fricke2020} and follow the shelf life of fresh produce \cite{haman2019}, and multi-temperature instruments give the temperature dependence of growth in one experiment \cite{wadso2011}. Three questions must be answered before the idea is used: how well one sample represents a lot, how stable the heat yield per cell is in milk, and how the detection count depends on the instrument and the sample volume. Each can be answered by a calorimetric experiment, and the generalized-time framework says which quantities such an experiment should report.

\subsection{Consistency with prior findings}
The sensitivity analyses that rank temperature first \cite{lau2022,buehler2018,qian2023,giannakourou2020} and our finding that biology dominates are two views of one system (Table~\ref{tab:findings}). When the temperature is unknown it dominates the spread of outcomes; once it is logged, Proposition~\ref{prop:clock} says the record has done its whole job, and what is left is biology. Our sensor sensitivity of \AnOneDeg\,\% per \textdegree C at 4\degC{} is steeper than the 10\,\% of a $Q_{10}=3$ model \cite{zou2023}, because the square-root law is steepest in the cold (Eq.~\eqref{eq:sens}); the two coincide near 10\degC.

\subsection{Hazard versus spoilage}
For spoilage flora a warm hour is expensive but the baseline rate is high; for \textit{Listeria}, growth at 4\degC{} is slow (generation time \NumLisGen\,h) but rises about \NumLisRatio-fold by 13\degC. The same benchmark and design rules hold for both, which supports treating temperature excursions as the common risk factor. This paper does not estimate illness risk, which depends on dose, virulence, host and the contamination event, none of which are modelled.

\subsection{Limitations}
(i)~The reference and the estimator both read the same generalized time (Proposition~\ref{prop:csl}), so the benchmark measures how much accuracy is lost to the sources listed, not whether a kinetic score models growth better; the structure-risk run relaxes one structural assumption only. (ii)~Strain spread, starting load, excursion statistics and the \textit{Listeria} lag are assumed. (iii)~Labels are constructed, with prevalence near 0.5 (spoilage) and \NumLisPrevMin--\NumLisPrevMax{} (\textit{Listeria}). (iv)~ComBase records come from heterogeneous studies and methods; maximum rates are not curves, so lag and responses to shifts are untested on them. (v)~One validation experiment reproduces fluctuating temperature, and no new plate counts were made. (vi)~Neither organism class represents spore-formers, which account for a large share of spoilage \cite{martin2021}. (vii)~For very small growth targets (the half-log \textit{Listeria} label) the estimator's structure, not biology alone, limits accuracy. (viii)~The calorimetric results are simulated: they assume a constant heat yield per cell, a fixed detection count and a sample that shares the biology of the trip it is used for, and no calorimetric measurement on milk was made.

\section{Conclusions}\label{sec:concl}
A temperature log can say much about how much growth the milk \emph{could} have had, and little about how much it \emph{did}. The mathematics says why: growth, and the heat it releases, depend on the temperature history only through a generalized time; a log measures that time almost perfectly even with a low-cost probe sampled hourly, while the rate scale and the budget against which it is compared vary from carton to carton. The simulation, anchored to published and ComBase data, says how much: unobserved biology costs \NumBaseGapMin--\NumBaseGapMax{} AUC; degree-hours come close to a growth model; sensor quality and sampling interval matter little; a probability is far better calibrated than an alarm; and an isothermal calorimetric detection time of the lot recovers \CalFracMin--\CalFracMax\,\% of what the log cannot see.

Five directions follow. First, Bayesian updating of the biological prior from occasional measurements (a plate count or a rapid adenosine triphosphate (ATP) test at the depot), for which the generalized-time reduction gives a one-dimensional likelihood. Second, a value-of-information analysis that asks which single measurement of biology most reduces misranking. Third, real logger data from Indian dairy routes, to replace the assumed excursion model. Fourth, calorimetric experiments on milk under stepwise temperature programmes, to test the master curve of Proposition~\ref{prop:calo} and to measure detection times of real lots. Fifth, challenge tests under fluctuating temperature with natural flora, to test the generalized-time reduction and the extra-lag hypothesis directly.

